\chapter{PyQCD 文档全量来源图谱}
\label{app:pyqcd-source-atlas}
\providecommand{\sourcepath}[1]{\seqsplit{#1}}
本附录把固定快照中的每一个文件映射到专著的来源单元和章节用途。它回答三个审计问题：文件是否存在且未被改写、知识内容进入哪一章、TeX/PDF 成对件与历史版本的重复关系是否被显式处理。This appendix is a machine-auditable ledger, not a bibliography by volume alone.
\begin{checkpointbox}{来源基线}
\small
\begin{itemize}
  \item 上游仓库：\url{https://gitee.com/zhangxin8069/PyQCD.git}。
  \item 固定提交：\texttt{0a2f85d30881fc8e}\allowbreak\texttt{ad60972c0afb9c02}\allowbreak\texttt{8abc348f}。
  \item 许可证：MIT，见 \sourcepath{docs/pyqcd\_reference/LICENSE.md}。
  \item 文件总数：166；内容文件：163；控制文件：3。
  \item 内容类型：77 个 TeX、80 个 PDF、6 个 Markdown；其中 76 对 TeX/PDF 成对保存，4 个 PDF 独立，1 个 TeX 独立。
  \item 文本规模：56\,850 行（TeX/Markdown）；PDF 规模：1\,566 页；原始字节：56\,165\,822。
  \item 完整性：163/163 个内容文件与吸收 manifest 的 SHA-256 一致；缺失文件 0，额外文件 0。
\end{itemize}
\end{checkpointbox}
\section{内容单元与章节映射}
成对 TeX/PDF 共享一个知识单元；PDF 仍保留为独立文件证据，并在逐文件账本中登记。每个单元在主批次中恰好出现一次。
\begin{description}[leftmargin=4.2em,labelwidth=3.7em,style=nextline,itemsep=0.35em]
\item[\texttt{U001}] \sourcepath{docs/pyqcd\_reference/PyQCD仓库结构与核心物理链解析.tex}
\par\noindent{\scriptsize
\textcolor{qcdsage}{伴生：}\sourcepath{docs/pyqcd\_reference/PyQCD仓库结构与核心物理链解析.pdf}。
类型：TeX + PDF；规模：166 行 / 4 页；目标：PyQCD 实现与复现。
}
\item[\texttt{U002}] \sourcepath{docs/pyqcd\_reference/absorb\_dev20\_gluon\_tmd\_audit\_20261022.md}
\par\noindent{\scriptsize
类型：Markdown；规模：378 行；目标：历史工作流与验证演进。
}
\item[\texttt{U003}] \sourcepath{docs/pyqcd\_reference/analy\_gpd\_20260928.tex}
\par\noindent{\scriptsize
\textcolor{qcdsage}{伴生：}\sourcepath{docs/pyqcd\_reference/analy\_gpd\_20260928.pdf}。
类型：TeX + PDF；规模：937 行 / 21 页；目标：PyQCD 实现与复现。
}
\item[\texttt{U004}] \sourcepath{docs/pyqcd\_reference/analy\_physics\_chain\_20260824.tex}
\par\noindent{\scriptsize
\textcolor{qcdsage}{伴生：}\sourcepath{docs/pyqcd\_reference/analy\_physics\_chain\_20260824.pdf}。
类型：TeX + PDF；规模：975 行 / 29 页；目标：PyQCD 实现与复现。
}
\item[\texttt{U005}] \sourcepath{docs/pyqcd\_reference/analy\_proton\_energy\_20260816.tex}
\par\noindent{\scriptsize
类型：独立 TeX；规模：300 行；目标：PyQCD 实现与复现。
}
\item[\texttt{U006}] \sourcepath{docs/pyqcd\_reference/analy\_pyqcd\_20260818.tex}
\par\noindent{\scriptsize
\textcolor{qcdsage}{伴生：}\sourcepath{docs/pyqcd\_reference/analy\_pyqcd\_20260818.pdf}。
类型：TeX + PDF；规模：324 行 / 11 页；目标：PyQCD 实现与复现。
}
\item[\texttt{U007}] \sourcepath{docs/pyqcd\_reference/analy\_torch\_parallel\_20260816.tex}
\par\noindent{\scriptsize
\textcolor{qcdsage}{伴生：}\sourcepath{docs/pyqcd\_reference/analy\_torch\_parallel\_20260816.pdf}。
类型：TeX + PDF；规模：442 行 / 9 页；目标：PyQCD 实现与复现。
}
\item[\texttt{U008}] \sourcepath{docs/pyqcd\_reference/from-logs/README.md}
\par\noindent{\scriptsize
类型：Markdown；规模：7 行；目标：学位论文综合与来源词库。
}
\item[\texttt{U009}] \sourcepath{docs/pyqcd\_reference/from-logs/logs/dev5/analy\_test9\_20260820.tex}
\par\noindent{\scriptsize
\textcolor{qcdsage}{伴生：}\sourcepath{docs/pyqcd\_reference/from-logs/logs/dev5/analy\_test9\_20260820.pdf}。
类型：TeX + PDF；规模：818 行 / 25 页；目标：历史工作流与验证演进。
}
\item[\texttt{U010}] \sourcepath{docs/pyqcd\_reference/from-logs/logs/dev5\_1/analy\_test9\_20260821.tex}
\par\noindent{\scriptsize
\textcolor{qcdsage}{伴生：}\sourcepath{docs/pyqcd\_reference/from-logs/logs/dev5\_1/analy\_test9\_20260821.pdf}。
类型：TeX + PDF；规模：863 行 / 33 页；目标：历史工作流与验证演进。
}
\item[\texttt{U011}] \sourcepath{docs/pyqcd\_reference/from-logs/logs/dev5\_2/analy\_test9\_20260822.tex}
\par\noindent{\scriptsize
\textcolor{qcdsage}{伴生：}\sourcepath{docs/pyqcd\_reference/from-logs/logs/dev5\_2/analy\_test9\_20260822.pdf}。
类型：TeX + PDF；规模：901 行 / 36 页；目标：历史工作流与验证演进。
}
\item[\texttt{U012}] \sourcepath{docs/pyqcd\_reference/from-logs/logs/donghx\_4150\_20260829/pdfcheck.md}
\par\noindent{\scriptsize
类型：Markdown；规模：70 行；目标：历史工作流与验证演进。
}
\item[\texttt{U013}] \sourcepath{docs/pyqcd\_reference/from-logs/logs/donghx\_4150\_20260829/summary.md}
\par\noindent{\scriptsize
类型：Markdown；规模：278 行；目标：历史工作流与验证演进。
}
\item[\texttt{U014}] \sourcepath{docs/pyqcd\_reference/from-logs/logs/stab0/stab0.tex}
\par\noindent{\scriptsize
\textcolor{qcdsage}{伴生：}\sourcepath{docs/pyqcd\_reference/from-logs/logs/stab0/stab0.pdf}。
类型：TeX + PDF；规模：589 行 / 13 页；目标：历史工作流与验证演进。
}
\item[\texttt{U015}] \sourcepath{docs/pyqcd\_reference/from-logs/logs/stab1/docs/stab1\_analysis.tex}
\par\noindent{\scriptsize
\textcolor{qcdsage}{伴生：}\sourcepath{docs/pyqcd\_reference/from-logs/logs/stab1/docs/stab1\_analysis.pdf}。
类型：TeX + PDF；规模：150 行 / 6 页；目标：历史工作流与验证演进。
}
\item[\texttt{U016}] \sourcepath{docs/pyqcd\_reference/from-logs/logs/stab1/v202608151614/stab1\_report.tex}
\par\noindent{\scriptsize
\textcolor{qcdsage}{伴生：}\sourcepath{docs/pyqcd\_reference/from-logs/logs/stab1/v202608151614/stab1\_report.pdf}。
类型：TeX + PDF；规模：74 行 / 6 页；目标：历史工作流与验证演进。
}
\item[\texttt{U017}] \sourcepath{docs/pyqcd\_reference/from-logs/logs/test6/docs/analy\_test6\_report\_20260816.tex}
\par\noindent{\scriptsize
\textcolor{qcdsage}{伴生：}\sourcepath{docs/pyqcd\_reference/from-logs/logs/test6/docs/analy\_test6\_report\_20260816.pdf}。
类型：TeX + PDF；规模：412 行 / 11 页；目标：真实数据复现与验证矩阵。
}
\item[\texttt{U018}] \sourcepath{docs/pyqcd\_reference/from-logs/logs/test7/v202608161007/test7\_report.tex}
\par\noindent{\scriptsize
\textcolor{qcdsage}{伴生：}\sourcepath{docs/pyqcd\_reference/from-logs/logs/test7/v202608161007/test7\_report.pdf}。
类型：TeX + PDF；规模：50 行 / 2 页；目标：真实数据复现与验证矩阵。
}
\item[\texttt{U019}] \sourcepath{docs/pyqcd\_reference/from-logs/logs/test8/v202608161117/test8\_report.tex}
\par\noindent{\scriptsize
\textcolor{qcdsage}{伴生：}\sourcepath{docs/pyqcd\_reference/from-logs/logs/test8/v202608161117/test8\_report.pdf}。
类型：TeX + PDF；规模：71 行 / 5 页；目标：真实数据复现与验证矩阵。
}
\item[\texttt{U020}] \sourcepath{docs/pyqcd\_reference/from-logs/logs/test9/test9\_analysis.tex}
\par\noindent{\scriptsize
\textcolor{qcdsage}{伴生：}\sourcepath{docs/pyqcd\_reference/from-logs/logs/test9/test9\_analysis.pdf}。
类型：TeX + PDF；规模：233 行 / 8 页；目标：真实数据复现与验证矩阵。
}
\item[\texttt{U021}] \sourcepath{docs/pyqcd\_reference/gevp\_convergence.pdf}
\par\noindent{\scriptsize
类型：独立 PDF；规模：1 页；目标：关联函数、谱学与有限体积。
}
\item[\texttt{U022}] \sourcepath{docs/pyqcd\_reference/gevp\_t0\_suppression.pdf}
\par\noindent{\scriptsize
类型：独立 PDF；规模：1 页；目标：关联函数、谱学与有限体积。
}
\item[\texttt{U023}] \sourcepath{docs/pyqcd\_reference/gradient\_flow.md}
\par\noindent{\scriptsize
类型：Markdown；规模：895 行；目标：梯度流、流时窗口与胶子 TMD。
}
\item[\texttt{U024}] \sourcepath{docs/pyqcd\_reference/pure\_pyqcd\_20260818.tex}
\par\noindent{\scriptsize
\textcolor{qcdsage}{伴生：}\sourcepath{docs/pyqcd\_reference/pure\_pyqcd\_20260818.pdf}。
类型：TeX + PDF；规模：304 行 / 10 页；目标：PyQCD 实现与复现。
}
\item[\texttt{U025}] \sourcepath{docs/pyqcd\_reference/report\_cmp1\_4150\_20260828.tex}
\par\noindent{\scriptsize
\textcolor{qcdsage}{伴生：}\sourcepath{docs/pyqcd\_reference/report\_cmp1\_4150\_20260828.pdf}。
类型：TeX + PDF；规模：507 行 / 16 页；目标：真实数据复现与验证矩阵。
}
\item[\texttt{U026}] \sourcepath{docs/pyqcd\_reference/report\_donghx\_4150\_reproduction\_20260829.tex}
\par\noindent{\scriptsize
\textcolor{qcdsage}{伴生：}\sourcepath{docs/pyqcd\_reference/report\_donghx\_4150\_reproduction\_20260829.pdf}。
类型：TeX + PDF；规模：617 行 / 23 页；目标：真实数据复现与验证矩阵。
}
\item[\texttt{U027}] \sourcepath{docs/pyqcd\_reference/report\_donghx\_4150\_reproduction\_20260830.tex}
\par\noindent{\scriptsize
\textcolor{qcdsage}{伴生：}\sourcepath{docs/pyqcd\_reference/report\_donghx\_4150\_reproduction\_20260830.pdf}。
类型：TeX + PDF；规模：732 行 / 27 页；目标：真实数据复现与验证矩阵。
}
\item[\texttt{U028}] \sourcepath{docs/pyqcd\_reference/report\_donghx\_pyqcd\_compare\_20260830.tex}
\par\noindent{\scriptsize
\textcolor{qcdsage}{伴生：}\sourcepath{docs/pyqcd\_reference/report\_donghx\_pyqcd\_compare\_20260830.pdf}。
类型：TeX + PDF；规模：361 行 / 12 页；目标：真实数据复现与验证矩阵。
}
\item[\texttt{U029}] \sourcepath{docs/pyqcd\_reference/report\_donghx\_pyqcd\_compare\_20260830\_2.tex}
\par\noindent{\scriptsize
\textcolor{qcdsage}{伴生：}\sourcepath{docs/pyqcd\_reference/report\_donghx\_pyqcd\_compare\_20260830\_2.pdf}。
类型：TeX + PDF；规模：351 行 / 10 页；目标：真实数据复现与验证矩阵。
}
\item[\texttt{U030}] \sourcepath{docs/pyqcd\_reference/report\_gluon\_tmd\_gradient\_flow\_20260828.tex}
\par\noindent{\scriptsize
\textcolor{qcdsage}{伴生：}\sourcepath{docs/pyqcd\_reference/report\_gluon\_tmd\_gradient\_flow\_20260828.pdf}。
类型：TeX + PDF；规模：789 行 / 22 页；目标：梯度流、流时窗口与胶子 TMD。
}
\item[\texttt{U031}] \sourcepath{docs/pyqcd\_reference/report\_gradient\_flow\_gluonpdf\_pyqcd\_20260909.tex}
\par\noindent{\scriptsize
\textcolor{qcdsage}{伴生：}\sourcepath{docs/pyqcd\_reference/report\_gradient\_flow\_gluonpdf\_pyqcd\_20260909.pdf}。
类型：TeX + PDF；规模：1029 行 / 17 页；目标：梯度流、流时窗口与胶子 TMD。
}
\item[\texttt{U032}] \sourcepath{docs/pyqcd\_reference/基于Gradient\_Flow的质子胶子PDF格点计算整体流程.md}
\par\noindent{\scriptsize
类型：Markdown；规模：661 行；目标：梯度流、流时窗口与胶子 TMD。
}
\item[\texttt{U033}] \sourcepath{docs/pyqcd\_reference/基于格点QCD对部分子分布函数的研究.pdf}
\par\noindent{\scriptsize
类型：独立 PDF；规模：56 页；目标：学位论文综合与来源词库。
}
\item[\texttt{U034}] \sourcepath{docs/pyqcd\_reference/基于格点QCD的横向动量依赖的部分子分布函数的研究.pdf}
\par\noindent{\scriptsize
类型：独立 PDF；规模：82 页；目标：学位论文综合与来源词库。
}
\item[\texttt{U035}] \sourcepath{docs/pyqcd\_reference/构造胶子准算符.tex}
\par\noindent{\scriptsize
\textcolor{qcdsage}{伴生：}\sourcepath{docs/pyqcd\_reference/构造胶子准算符.pdf}。
类型：TeX + PDF；规模：562 行 / 14 页；目标：胶子算符、TMD 与自旋。
}
\item[\texttt{U036}] \sourcepath{docs/pyqcd\_reference/格点QCD中的3pt构造.tex}
\par\noindent{\scriptsize
\textcolor{qcdsage}{伴生：}\sourcepath{docs/pyqcd\_reference/格点QCD中的3pt构造.pdf}。
类型：TeX + PDF；规模：783 行 / 20 页；目标：关联函数、谱学与有限体积。
}
\item[\texttt{U037}] \sourcepath{docs/pyqcd\_reference/格点QCD中的DGLAP演化方程.tex}
\par\noindent{\scriptsize
\textcolor{qcdsage}{伴生：}\sourcepath{docs/pyqcd\_reference/格点QCD中的DGLAP演化方程.pdf}。
类型：TeX + PDF；规模：688 行 / 18 页；目标：PDF、OPE、LaMET 与 GPD。
}
\item[\texttt{U038}] \sourcepath{docs/pyqcd\_reference/格点QCD中的GEVP.tex}
\par\noindent{\scriptsize
\textcolor{qcdsage}{伴生：}\sourcepath{docs/pyqcd\_reference/格点QCD中的GEVP.pdf}。
类型：TeX + PDF；规模：924 行 / 24 页；目标：关联函数、谱学与有限体积。
}
\item[\texttt{U039}] \sourcepath{docs/pyqcd\_reference/格点QCD中的OPE算符.tex}
\par\noindent{\scriptsize
\textcolor{qcdsage}{伴生：}\sourcepath{docs/pyqcd\_reference/格点QCD中的OPE算符.pdf}。
类型：TeX + PDF；规模：526 行 / 14 页；目标：PDF、OPE、LaMET 与 GPD。
}
\item[\texttt{U040}] \sourcepath{docs/pyqcd\_reference/格点QCD中的Symanzik有效理论.tex}
\par\noindent{\scriptsize
\textcolor{qcdsage}{伴生：}\sourcepath{docs/pyqcd\_reference/格点QCD中的Symanzik有效理论.pdf}。
类型：TeX + PDF；规模：449 行 / 13 页；目标：格点规范理论基础与系统误差。
}
\item[\texttt{U041}] \sourcepath{docs/pyqcd\_reference/格点QCD中的TMD\_PDF.tex}
\par\noindent{\scriptsize
\textcolor{qcdsage}{伴生：}\sourcepath{docs/pyqcd\_reference/格点QCD中的TMD\_PDF.pdf}。
类型：TeX + PDF；规模：937 行 / 30 页；目标：胶子算符、TMD 与自旋。
}
\item[\texttt{U042}] \sourcepath{docs/pyqcd\_reference/格点QCD中的UV涨落.tex}
\par\noindent{\scriptsize
\textcolor{qcdsage}{伴生：}\sourcepath{docs/pyqcd\_reference/格点QCD中的UV涨落.pdf}。
类型：TeX + PDF；规模：472 行 / 13 页；目标：格点规范理论基础与系统误差。
}
\item[\texttt{U043}] \sourcepath{docs/pyqcd\_reference/格点QCD中的Wilson线.tex}
\par\noindent{\scriptsize
\textcolor{qcdsage}{伴生：}\sourcepath{docs/pyqcd\_reference/格点QCD中的Wilson线.pdf}。
类型：TeX + PDF；规模：712 行 / 21 页；目标：格点规范理论基础与系统误差。
}
\item[\texttt{U044}] \sourcepath{docs/pyqcd\_reference/格点QCD中的luscher参数化.tex}
\par\noindent{\scriptsize
\textcolor{qcdsage}{伴生：}\sourcepath{docs/pyqcd\_reference/格点QCD中的luscher参数化.pdf}。
类型：TeX + PDF；规模：603 行 / 17 页；目标：关联函数、谱学与有限体积。
}
\item[\texttt{U045}] \sourcepath{docs/pyqcd\_reference/格点QCD中的smear算法.tex}
\par\noindent{\scriptsize
\textcolor{qcdsage}{伴生：}\sourcepath{docs/pyqcd\_reference/格点QCD中的smear算法.pdf}。
类型：TeX + PDF；规模：630 行 / 16 页；目标：数值算法、统计推断与外推。
}
\item[\texttt{U046}] \sourcepath{docs/pyqcd\_reference/格点QCD中的传播子求解.tex}
\par\noindent{\scriptsize
\textcolor{qcdsage}{伴生：}\sourcepath{docs/pyqcd\_reference/格点QCD中的传播子求解.pdf}。
类型：TeX + PDF；规模：269 行 / 10 页；目标：数值算法、统计推断与外推。
}
\item[\texttt{U047}] \sourcepath{docs/pyqcd\_reference/格点QCD中的光锥PDF、准PDF、赝PDF.tex}
\par\noindent{\scriptsize
\textcolor{qcdsage}{伴生：}\sourcepath{docs/pyqcd\_reference/格点QCD中的光锥PDF、准PDF、赝PDF.pdf}。
类型：TeX + PDF；规模：705 行 / 20 页；目标：PDF、OPE、LaMET 与 GPD。
}
\item[\texttt{U048}] \sourcepath{docs/pyqcd\_reference/格点QCD中的光锥PDF与quasi\_PDF.tex}
\par\noindent{\scriptsize
\textcolor{qcdsage}{伴生：}\sourcepath{docs/pyqcd\_reference/格点QCD中的光锥PDF与quasi\_PDF.pdf}。
类型：TeX + PDF；规模：413 行 / 12 页；目标：PDF、OPE、LaMET 与 GPD。
}
\item[\texttt{U049}] \sourcepath{docs/pyqcd\_reference/格点QCD中的光锥与光前.tex}
\par\noindent{\scriptsize
\textcolor{qcdsage}{伴生：}\sourcepath{docs/pyqcd\_reference/格点QCD中的光锥与光前.pdf}。
类型：TeX + PDF；规模：783 行 / 18 页；目标：PDF、OPE、LaMET 与 GPD。
}
\item[\texttt{U050}] \sourcepath{docs/pyqcd\_reference/格点QCD中的关联函数.tex}
\par\noindent{\scriptsize
\textcolor{qcdsage}{伴生：}\sourcepath{docs/pyqcd\_reference/格点QCD中的关联函数.pdf}。
类型：TeX + PDF；规模：604 行 / 15 页；目标：关联函数、谱学与有限体积。
}
\item[\texttt{U051}] \sourcepath{docs/pyqcd\_reference/格点QCD中的场强张量.tex}
\par\noindent{\scriptsize
\textcolor{qcdsage}{伴生：}\sourcepath{docs/pyqcd\_reference/格点QCD中的场强张量.pdf}。
类型：TeX + PDF；规模：576 行 / 14 页；目标：格点规范理论基础与系统误差。
}
\item[\texttt{U052}] \sourcepath{docs/pyqcd\_reference/格点QCD中的外推.tex}
\par\noindent{\scriptsize
\textcolor{qcdsage}{伴生：}\sourcepath{docs/pyqcd\_reference/格点QCD中的外推.pdf}。
类型：TeX + PDF；规模：567 行 / 16 页；目标：数值算法、统计推断与外推。
}
\item[\texttt{U053}] \sourcepath{docs/pyqcd\_reference/格点QCD中的大动量有效理论.tex}
\par\noindent{\scriptsize
\textcolor{qcdsage}{伴生：}\sourcepath{docs/pyqcd\_reference/格点QCD中的大动量有效理论.pdf}。
类型：TeX + PDF；规模：498 行 / 13 页；目标：PDF、OPE、LaMET 与 GPD。
}
\item[\texttt{U054}] \sourcepath{docs/pyqcd\_reference/格点QCD中的夸克算符.tex}
\par\noindent{\scriptsize
\textcolor{qcdsage}{伴生：}\sourcepath{docs/pyqcd\_reference/格点QCD中的夸克算符.pdf}。
类型：TeX + PDF；规模：457 行 / 13 页；目标：PDF、OPE、LaMET 与 GPD。
}
\item[\texttt{U055}] \sourcepath{docs/pyqcd\_reference/格点QCD中的常见名词短语与缩写.tex}
\par\noindent{\scriptsize
\textcolor{qcdsage}{伴生：}\sourcepath{docs/pyqcd\_reference/格点QCD中的常见名词短语与缩写.pdf}。
类型：TeX + PDF；规模：414 行 / 26 页；目标：学位论文综合与来源词库。
}
\item[\texttt{U056}] \sourcepath{docs/pyqcd\_reference/格点QCD中的常见费曼图.tex}
\par\noindent{\scriptsize
\textcolor{qcdsage}{伴生：}\sourcepath{docs/pyqcd\_reference/格点QCD中的常见费曼图.pdf}。
类型：TeX + PDF；规模：495 行 / 10 页；目标：关联函数、谱学与有限体积。
}
\item[\texttt{U057}] \sourcepath{docs/pyqcd\_reference/格点QCD中的广义部分子分布.tex}
\par\noindent{\scriptsize
\textcolor{qcdsage}{伴生：}\sourcepath{docs/pyqcd\_reference/格点QCD中的广义部分子分布.pdf}。
类型：TeX + PDF；规模：1354 行 / 25 页；目标：PDF、OPE、LaMET 与 GPD。
}
\item[\texttt{U058}] \sourcepath{docs/pyqcd\_reference/格点QCD中的强子谱学.tex}
\par\noindent{\scriptsize
\textcolor{qcdsage}{伴生：}\sourcepath{docs/pyqcd\_reference/格点QCD中的强子谱学.pdf}。
类型：TeX + PDF；规模：353 行 / 11 页；目标：关联函数、谱学与有限体积。
}
\item[\texttt{U059}] \sourcepath{docs/pyqcd\_reference/格点QCD中的形状因子.tex}
\par\noindent{\scriptsize
\textcolor{qcdsage}{伴生：}\sourcepath{docs/pyqcd\_reference/格点QCD中的形状因子.pdf}。
类型：TeX + PDF；规模：499 行 / 14 页；目标：关联函数、谱学与有限体积。
}
\item[\texttt{U060}] \sourcepath{docs/pyqcd\_reference/格点QCD中的标度参数.tex}
\par\noindent{\scriptsize
\textcolor{qcdsage}{伴生：}\sourcepath{docs/pyqcd\_reference/格点QCD中的标度参数.pdf}。
类型：TeX + PDF；规模：456 行 / 13 页；目标：格点规范理论基础与系统误差。
}
\item[\texttt{U061}] \sourcepath{docs/pyqcd\_reference/格点QCD中的梯度流重整化.tex}
\par\noindent{\scriptsize
\textcolor{qcdsage}{伴生：}\sourcepath{docs/pyqcd\_reference/格点QCD中的梯度流重整化.pdf}。
类型：TeX + PDF；规模：626 行 / 15 页；目标：梯度流、流时窗口与胶子 TMD。
}
\item[\texttt{U062}] \sourcepath{docs/pyqcd\_reference/格点QCD中的维克收缩与傅里叶变换.tex}
\par\noindent{\scriptsize
\textcolor{qcdsage}{伴生：}\sourcepath{docs/pyqcd\_reference/格点QCD中的维克收缩与傅里叶变换.pdf}。
类型：TeX + PDF；规模：469 行 / 13 页；目标：关联函数、谱学与有限体积。
}
\item[\texttt{U063}] \sourcepath{docs/pyqcd\_reference/格点QCD中的胶子极化.tex}
\par\noindent{\scriptsize
\textcolor{qcdsage}{伴生：}\sourcepath{docs/pyqcd\_reference/格点QCD中的胶子极化.pdf}。
类型：TeX + PDF；规模：694 行 / 16 页；目标：胶子算符、TMD 与自旋。
}
\item[\texttt{U064}] \sourcepath{docs/pyqcd\_reference/格点QCD中的胶子算符.tex}
\par\noindent{\scriptsize
\textcolor{qcdsage}{伴生：}\sourcepath{docs/pyqcd\_reference/格点QCD中的胶子算符.pdf}。
类型：TeX + PDF；规模：618 行 / 15 页；目标：胶子算符、TMD 与自旋。
}
\item[\texttt{U065}] \sourcepath{docs/pyqcd\_reference/格点QCD中的色散关系.tex}
\par\noindent{\scriptsize
\textcolor{qcdsage}{伴生：}\sourcepath{docs/pyqcd\_reference/格点QCD中的色散关系.pdf}。
类型：TeX + PDF；规模：496 行 / 15 页；目标：格点规范理论基础与系统误差。
}
\item[\texttt{U066}] \sourcepath{docs/pyqcd\_reference/格点QCD中的蒙卡方法.tex}
\par\noindent{\scriptsize
\textcolor{qcdsage}{伴生：}\sourcepath{docs/pyqcd\_reference/格点QCD中的蒙卡方法.pdf}。
类型：TeX + PDF；规模：347 行 / 11 页；目标：格点规范理论基础与系统误差。
}
\item[\texttt{U067}] \sourcepath{docs/pyqcd\_reference/格点QCD中的衰变道.tex}
\par\noindent{\scriptsize
\textcolor{qcdsage}{伴生：}\sourcepath{docs/pyqcd\_reference/格点QCD中的衰变道.pdf}。
类型：TeX + PDF；规模：474 行 / 15 页；目标：关联函数、谱学与有限体积。
}
\item[\texttt{U068}] \sourcepath{docs/pyqcd\_reference/格点QCD中的误差统计.tex}
\par\noindent{\scriptsize
\textcolor{qcdsage}{伴生：}\sourcepath{docs/pyqcd\_reference/格点QCD中的误差统计.pdf}。
类型：TeX + PDF；规模：472 行 / 15 页；目标：数值算法、统计推断与外推。
}
\item[\texttt{U069}] \sourcepath{docs/pyqcd\_reference/格点QCD中的费米子方案.tex}
\par\noindent{\scriptsize
\textcolor{qcdsage}{伴生：}\sourcepath{docs/pyqcd\_reference/格点QCD中的费米子方案.pdf}。
类型：TeX + PDF；规模：999 行 / 24 页；目标：格点规范理论基础与系统误差。
}
\item[\texttt{U070}] \sourcepath{docs/pyqcd\_reference/格点QCD中的连通图与非连通图.tex}
\par\noindent{\scriptsize
\textcolor{qcdsage}{伴生：}\sourcepath{docs/pyqcd\_reference/格点QCD中的连通图与非连通图.pdf}。
类型：TeX + PDF；规模：531 行 / 16 页；目标：关联函数、谱学与有限体积。
}
\item[\texttt{U071}] \sourcepath{docs/pyqcd\_reference/格点QCD中的部分子分布函数.tex}
\par\noindent{\scriptsize
\textcolor{qcdsage}{伴生：}\sourcepath{docs/pyqcd\_reference/格点QCD中的部分子分布函数.pdf}。
类型：TeX + PDF；规模：576 行 / 16 页；目标：PDF、OPE、LaMET 与 GPD。
}
\item[\texttt{U072}] \sourcepath{docs/pyqcd\_reference/格点QCD中的重夸克有效理论.tex}
\par\noindent{\scriptsize
\textcolor{qcdsage}{伴生：}\sourcepath{docs/pyqcd\_reference/格点QCD中的重夸克有效理论.pdf}。
类型：TeX + PDF；规模：426 行 / 12 页；目标：格点规范理论基础与系统误差。
}
\item[\texttt{U073}] \sourcepath{docs/pyqcd\_reference/格点QCD中的重整化.tex}
\par\noindent{\scriptsize
\textcolor{qcdsage}{伴生：}\sourcepath{docs/pyqcd\_reference/格点QCD中的重整化.pdf}。
类型：TeX + PDF；规模：741 行 / 21 页；目标：格点规范理论基础与系统误差。
}
\item[\texttt{U074}] \sourcepath{docs/pyqcd\_reference/格点QCD中的重采样方法.tex}
\par\noindent{\scriptsize
\textcolor{qcdsage}{伴生：}\sourcepath{docs/pyqcd\_reference/格点QCD中的重采样方法.pdf}。
类型：TeX + PDF；规模：507 行 / 16 页；目标：数值算法、统计推断与外推。
}
\item[\texttt{U075}] \sourcepath{docs/pyqcd\_reference/格点QCD蒸馏方法解析.tex}
\par\noindent{\scriptsize
\textcolor{qcdsage}{伴生：}\sourcepath{docs/pyqcd\_reference/格点QCD蒸馏方法解析.pdf}。
类型：TeX + PDF；规模：586 行 / 14 页；目标：关联函数、谱学与有限体积。
}
\item[\texttt{U076}] \sourcepath{docs/pyqcd\_reference/格点上计算胶子准算符.tex}
\par\noindent{\scriptsize
\textcolor{qcdsage}{伴生：}\sourcepath{docs/pyqcd\_reference/格点上计算胶子准算符.pdf}。
类型：TeX + PDF；规模：947 行 / 22 页；目标：胶子算符、TMD 与自旋。
}
\item[\texttt{U077}] \sourcepath{docs/pyqcd\_reference/理论解析与工作流-v20260911.tex}
\par\noindent{\scriptsize
\textcolor{qcdsage}{伴生：}\sourcepath{docs/pyqcd\_reference/理论解析与工作流-v20260911.pdf}。
类型：TeX + PDF；规模：5957 行 / 100 页；目标：学位论文综合与来源词库。
}
\item[\texttt{U078}] \sourcepath{docs/pyqcd\_reference/理论解析与工作流.tex}
\par\noindent{\scriptsize
\textcolor{qcdsage}{伴生：}\sourcepath{docs/pyqcd\_reference/理论解析与工作流.pdf}。
类型：TeX + PDF；规模：3837 行 / 89 页；目标：学位论文综合与来源词库。
}
\item[\texttt{U079}] \sourcepath{docs/pyqcd\_reference/胶子PDF完整推导.tex}
\par\noindent{\scriptsize
\textcolor{qcdsage}{伴生：}\sourcepath{docs/pyqcd\_reference/胶子PDF完整推导.pdf}。
类型：TeX + PDF；规模：742 行 / 15 页；目标：PDF、OPE、LaMET 与 GPD。
}
\item[\texttt{U080}] \sourcepath{docs/pyqcd\_reference/胶子PDF幻灯片.tex}
\par\noindent{\scriptsize
\textcolor{qcdsage}{伴生：}\sourcepath{docs/pyqcd\_reference/胶子PDF幻灯片.pdf}。
类型：TeX + PDF；规模：357 行 / 10 页；目标：胶子算符、TMD 与自旋。
}
\item[\texttt{U081}] \sourcepath{docs/pyqcd\_reference/胶子PDF笔记.tex}
\par\noindent{\scriptsize
\textcolor{qcdsage}{伴生：}\sourcepath{docs/pyqcd\_reference/胶子PDF笔记.pdf}。
类型：TeX + PDF；规模：445 行 / 9 页；目标：PDF、OPE、LaMET 与 GPD。
}
\item[\texttt{U082}] \sourcepath{docs/pyqcd\_reference/胶子PDF计算代码解析.tex}
\par\noindent{\scriptsize
\textcolor{qcdsage}{伴生：}\sourcepath{docs/pyqcd\_reference/胶子PDF计算代码解析.pdf}。
类型：TeX + PDF；规模：1218 行 / 30 页；目标：PyQCD 实现与复现。
}
\item[\texttt{U083}] \sourcepath{docs/pyqcd\_reference/胶子PDF连续极限.tex}
\par\noindent{\scriptsize
\textcolor{qcdsage}{伴生：}\sourcepath{docs/pyqcd\_reference/胶子PDF连续极限.pdf}。
类型：TeX + PDF；规模：578 行 / 16 页；目标：胶子算符、TMD 与自旋。
}
\item[\texttt{U084}] \sourcepath{docs/pyqcd\_reference/胶子PDF连续极限幻灯片.tex}
\par\noindent{\scriptsize
\textcolor{qcdsage}{伴生：}\sourcepath{docs/pyqcd\_reference/胶子PDF连续极限幻灯片.pdf}。
类型：TeX + PDF；规模：1576 行 / 50 页；目标：胶子算符、TMD 与自旋。
}
\item[\texttt{U085}] \sourcepath{docs/pyqcd\_reference/胶子PDF连续极限论文.tex}
\par\noindent{\scriptsize
\textcolor{qcdsage}{伴生：}\sourcepath{docs/pyqcd\_reference/胶子PDF连续极限论文.pdf}。
类型：TeX + PDF；规模：573 行 / 16 页；目标：胶子算符、TMD 与自旋。
}
\item[\texttt{U086}] \sourcepath{docs/pyqcd\_reference/胶子TMD梯度流连续极限.tex}
\par\noindent{\scriptsize
\textcolor{qcdsage}{伴生：}\sourcepath{docs/pyqcd\_reference/胶子TMD梯度流连续极限.pdf}。
类型：TeX + PDF；规模：1481 行 / 35 页；目标：胶子算符、TMD 与自旋。
}
\item[\texttt{U087}] \sourcepath{docs/pyqcd\_reference/质子自旋危机解析.tex}
\par\noindent{\scriptsize
\textcolor{qcdsage}{伴生：}\sourcepath{docs/pyqcd\_reference/质子自旋危机解析.pdf}。
类型：TeX + PDF；规模：564 行 / 14 页；目标：胶子算符、TMD 与自旋。
}
\end{description}
\section{逐文件哈希与规模账本}
本账本以文件而不是知识单元为粒度，确保 PDF 见证件、历史版本和 Markdown 工作记录没有被去重后从审计记录中消失。
\begin{description}[leftmargin=4.2em,labelwidth=3.7em,style=nextline,itemsep=0.18em]
\item[\texttt{001}] \sourcepath{docs/pyqcd\_reference/PyQCD仓库结构与核心物理链解析.pdf}
\par\noindent{\scriptsize PDF 见证；规模：4 页；字节：206\,963；SHA-256：\texttt{2d09cac620d75bf8}\allowbreak\texttt{7b01967fe155f35e}\allowbreak\texttt{65906959656fbf1a}\allowbreak\texttt{f8693cec45cc3318}。}
\item[\texttt{002}] \sourcepath{docs/pyqcd\_reference/PyQCD仓库结构与核心物理链解析.tex}
\par\noindent{\scriptsize TeX 正文；规模：166 行；字节：8\,336；SHA-256：\texttt{134ce0210d6793fc}\allowbreak\texttt{feda59b9fce7fdb4}\allowbreak\texttt{797a5031e300fd4d}\allowbreak\texttt{6905c116e960eef6}。}
\item[\texttt{003}] \sourcepath{docs/pyqcd\_reference/absorb\_dev20\_gluon\_tmd\_audit\_20261022.md}
\par\noindent{\scriptsize Markdown；规模：378 行；字节：18\,531；SHA-256：\texttt{4ee378fef6d34d8a}\allowbreak\texttt{54d91ded0172113c}\allowbreak\texttt{0b0e18d5ccbc02fb}\allowbreak\texttt{35b4049162c550ca}。}
\item[\texttt{004}] \sourcepath{docs/pyqcd\_reference/analy\_gpd\_20260928.pdf}
\par\noindent{\scriptsize PDF 见证；规模：21 页；字节：635\,062；SHA-256：\texttt{2b4c554ef47f5022}\allowbreak\texttt{f8f710058abd0f2f}\allowbreak\texttt{9a82d24dd5097613}\allowbreak\texttt{cf03b5e35d8dd14e}。}
\item[\texttt{005}] \sourcepath{docs/pyqcd\_reference/analy\_gpd\_20260928.tex}
\par\noindent{\scriptsize TeX 正文；规模：937 行；字节：45\,915；SHA-256：\texttt{ad2e68e844b6cf05}\allowbreak\texttt{26e0c2253ed0fb38}\allowbreak\texttt{0818763d6d576d00}\allowbreak\texttt{9719a6120799e5d1}。}
\item[\texttt{006}] \sourcepath{docs/pyqcd\_reference/analy\_physics\_chain\_20260824.pdf}
\par\noindent{\scriptsize PDF 见证；规模：29 页；字节：808\,442；SHA-256：\texttt{5a55abe4a7f4e1b0}\allowbreak\texttt{4fc26b934d7a68e6}\allowbreak\texttt{f81d564c704c080d}\allowbreak\texttt{e14d48e245cc2b71}。}
\item[\texttt{007}] \sourcepath{docs/pyqcd\_reference/analy\_physics\_chain\_20260824.tex}
\par\noindent{\scriptsize TeX 正文；规模：975 行；字节：56\,818；SHA-256：\texttt{9e5e3a1ff9ae08c5}\allowbreak\texttt{93222bc7827953eb}\allowbreak\texttt{23455a5f0f81e653}\allowbreak\texttt{65d1439824c7362b}。}
\item[\texttt{008}] \sourcepath{docs/pyqcd\_reference/analy\_proton\_energy\_20260816.tex}
\par\noindent{\scriptsize 独立 TeX；规模：300 行；字节：17\,301；SHA-256：\texttt{1c1be3bbfa90e622}\allowbreak\texttt{75b3f5d90a5d6ab1}\allowbreak\texttt{0ea7f71d2951dfb2}\allowbreak\texttt{068489135e7527d6}。}
\item[\texttt{009}] \sourcepath{docs/pyqcd\_reference/analy\_pyqcd\_20260818.pdf}
\par\noindent{\scriptsize PDF 见证；规模：11 页；字节：379\,821；SHA-256：\texttt{9b0dbc554d95d51e}\allowbreak\texttt{f24c90980c96e17c}\allowbreak\texttt{bd2c62ca7f4e7132}\allowbreak\texttt{cba1c6dd25b2526a}。}
\item[\texttt{010}] \sourcepath{docs/pyqcd\_reference/analy\_pyqcd\_20260818.tex}
\par\noindent{\scriptsize TeX 正文；规模：324 行；字节：18\,759；SHA-256：\texttt{fdacbc4e78a71fc9}\allowbreak\texttt{739a894fd69f7519}\allowbreak\texttt{186ce5504c92f9be}\allowbreak\texttt{ed10d082bfacd383}。}
\item[\texttt{011}] \sourcepath{docs/pyqcd\_reference/analy\_torch\_parallel\_20260816.pdf}
\par\noindent{\scriptsize PDF 见证；规模：9 页；字节：318\,869；SHA-256：\texttt{b88ab30252bf826f}\allowbreak\texttt{ea418721e4eeadf5}\allowbreak\texttt{176b1e22a910ec30}\allowbreak\texttt{7dde07ebbc6757c3}。}
\item[\texttt{012}] \sourcepath{docs/pyqcd\_reference/analy\_torch\_parallel\_20260816.tex}
\par\noindent{\scriptsize TeX 正文；规模：442 行；字节：21\,987；SHA-256：\texttt{cae2a64f7616b320}\allowbreak\texttt{cf4f213d8d938e37}\allowbreak\texttt{c48971f7ade562d8}\allowbreak\texttt{10126b79158340c2}。}
\item[\texttt{013}] \sourcepath{docs/pyqcd\_reference/from-logs/README.md}
\par\noindent{\scriptsize Markdown；规模：7 行；字节：327；SHA-256：\texttt{fc0b384080783923}\allowbreak\texttt{644288eb7864dcbf}\allowbreak\texttt{416cd9493f76e430}\allowbreak\texttt{bf11c125a1a08bd2}。}
\item[\texttt{014}] \sourcepath{docs/pyqcd\_reference/from-logs/logs/dev5/analy\_test9\_20260820.pdf}
\par\noindent{\scriptsize PDF 见证；规模：25 页；字节：1\,202\,601；SHA-256：\texttt{9e5d8f06527bace0}\allowbreak\texttt{2adc3195ab3a4a71}\allowbreak\texttt{4822e6183330abe3}\allowbreak\texttt{984d345786e9a534}。}
\item[\texttt{015}] \sourcepath{docs/pyqcd\_reference/from-logs/logs/dev5/analy\_test9\_20260820.tex}
\par\noindent{\scriptsize TeX 正文；规模：818 行；字节：64\,359；SHA-256：\texttt{52bd8bf496a7a4b0}\allowbreak\texttt{9738c1ec5259c265}\allowbreak\texttt{6712fca9b7746f8f}\allowbreak\texttt{e490bec1589d93ca}。}
\item[\texttt{016}] \sourcepath{docs/pyqcd\_reference/from-logs/logs/dev5\_1/analy\_test9\_20260821.pdf}
\par\noindent{\scriptsize PDF 见证；规模：33 页；字节：1\,245\,025；SHA-256：\texttt{a1572ea6320000a6}\allowbreak\texttt{ff8a6a6254078f52}\allowbreak\texttt{d89df71bc0c96453}\allowbreak\texttt{a8d0ae4989461572}。}
\item[\texttt{017}] \sourcepath{docs/pyqcd\_reference/from-logs/logs/dev5\_1/analy\_test9\_20260821.tex}
\par\noindent{\scriptsize TeX 正文；规模：863 行；字节：69\,355；SHA-256：\texttt{8bcb6539891ea2a2}\allowbreak\texttt{ae7e457520483476}\allowbreak\texttt{62ed1f96b7dd0ac0}\allowbreak\texttt{9bda309b145bba59}。}
\item[\texttt{018}] \sourcepath{docs/pyqcd\_reference/from-logs/logs/dev5\_2/analy\_test9\_20260822.pdf}
\par\noindent{\scriptsize PDF 见证；规模：36 页；字节：1\,280\,233；SHA-256：\texttt{d15cb8be273601f7}\allowbreak\texttt{129f99cbcba19e25}\allowbreak\texttt{988e81ce11b760ef}\allowbreak\texttt{c12785fff2b845cc}。}
\item[\texttt{019}] \sourcepath{docs/pyqcd\_reference/from-logs/logs/dev5\_2/analy\_test9\_20260822.tex}
\par\noindent{\scriptsize TeX 正文；规模：901 行；字节：74\,317；SHA-256：\texttt{8d0dc42c9031ed40}\allowbreak\texttt{3e2c4c739bb0a135}\allowbreak\texttt{7a03a4f528090fb5}\allowbreak\texttt{b2a1bc2f34b6d306}。}
\item[\texttt{020}] \sourcepath{docs/pyqcd\_reference/from-logs/logs/donghx\_4150\_20260829/pdfcheck.md}
\par\noindent{\scriptsize Markdown；规模：70 行；字节：2\,517；SHA-256：\texttt{c084c01cb5831f69}\allowbreak\texttt{dff5a0bed62685a9}\allowbreak\texttt{15c610d735b53371}\allowbreak\texttt{7679b2304e94595e}。}
\item[\texttt{021}] \sourcepath{docs/pyqcd\_reference/from-logs/logs/donghx\_4150\_20260829/summary.md}
\par\noindent{\scriptsize Markdown；规模：278 行；字节：14\,427；SHA-256：\texttt{1a21af7865602318}\allowbreak\texttt{f86c2b296f668bf7}\allowbreak\texttt{6b36b77876defb17}\allowbreak\texttt{341b34d284eda38b}。}
\item[\texttt{022}] \sourcepath{docs/pyqcd\_reference/from-logs/logs/stab0/stab0.pdf}
\par\noindent{\scriptsize PDF 见证；规模：13 页；字节：2\,463\,407；SHA-256：\texttt{1c82b836b99db690}\allowbreak\texttt{685e481b61fcaff6}\allowbreak\texttt{6138ae68d5c4cc62}\allowbreak\texttt{f9d8309af9d3dc7a}。}
\item[\texttt{023}] \sourcepath{docs/pyqcd\_reference/from-logs/logs/stab0/stab0.tex}
\par\noindent{\scriptsize TeX 正文；规模：589 行；字节：28\,489；SHA-256：\texttt{1b127b5b56560cf1}\allowbreak\texttt{d355f8d4fff0c4fd}\allowbreak\texttt{fe3c79498de0ab17}\allowbreak\texttt{280b489d7478fec0}。}
\item[\texttt{024}] \sourcepath{docs/pyqcd\_reference/from-logs/logs/stab1/docs/stab1\_analysis.pdf}
\par\noindent{\scriptsize PDF 见证；规模：6 页；字节：497\,785；SHA-256：\texttt{fb98c1ef65a170e9}\allowbreak\texttt{362f04bb991d86a9}\allowbreak\texttt{70eb68c1352160ef}\allowbreak\texttt{bcc3a5fcd6c04eec}。}
\item[\texttt{025}] \sourcepath{docs/pyqcd\_reference/from-logs/logs/stab1/docs/stab1\_analysis.tex}
\par\noindent{\scriptsize TeX 正文；规模：150 行；字节：7\,900；SHA-256：\texttt{08914a4fd86d1c6c}\allowbreak\texttt{8c425ac074846344}\allowbreak\texttt{61ae2ea8236813a8}\allowbreak\texttt{31f85df07066602f}。}
\item[\texttt{026}] \sourcepath{docs/pyqcd\_reference/from-logs/logs/stab1/v202608151614/stab1\_report.pdf}
\par\noindent{\scriptsize PDF 见证；规模：6 页；字节：538\,176；SHA-256：\texttt{7d5a713e6a56c8f3}\allowbreak\texttt{f4e02b94e59ad3ef}\allowbreak\texttt{a5fb4e6498665865}\allowbreak\texttt{eead8d964af53281}。}
\item[\texttt{027}] \sourcepath{docs/pyqcd\_reference/from-logs/logs/stab1/v202608151614/stab1\_report.tex}
\par\noindent{\scriptsize TeX 正文；规模：74 行；字节：2\,855；SHA-256：\texttt{3bc2467fa08b5f42}\allowbreak\texttt{1e5e3a2fb072306a}\allowbreak\texttt{984ccf0fe2687fe6}\allowbreak\texttt{d7f12aa1f30eeef9}。}
\item[\texttt{028}] \sourcepath{docs/pyqcd\_reference/from-logs/logs/test6/docs/analy\_test6\_report\_20260816.pdf}
\par\noindent{\scriptsize PDF 见证；规模：11 页；字节：674\,241；SHA-256：\texttt{a40e002a496987b6}\allowbreak\texttt{f7e1951a55cc975e}\allowbreak\texttt{9fba74e731585deb}\allowbreak\texttt{1f71a49ddc5c5925}。}
\item[\texttt{029}] \sourcepath{docs/pyqcd\_reference/from-logs/logs/test6/docs/analy\_test6\_report\_20260816.tex}
\par\noindent{\scriptsize TeX 正文；规模：412 行；字节：20\,924；SHA-256：\texttt{a36d4b185f94cc4c}\allowbreak\texttt{3b32b27472bc3bf3}\allowbreak\texttt{92d3abf531d43540}\allowbreak\texttt{a02ef7d01d175172}。}
\item[\texttt{030}] \sourcepath{docs/pyqcd\_reference/from-logs/logs/test7/v202608161007/test7\_report.pdf}
\par\noindent{\scriptsize PDF 见证；规模：2 页；字节：78\,209；SHA-256：\texttt{506531296ca5f609}\allowbreak\texttt{1118dfa9e9c53c29}\allowbreak\texttt{75e76ef019edb651}\allowbreak\texttt{d274ff0630c0beb6}。}
\item[\texttt{031}] \sourcepath{docs/pyqcd\_reference/from-logs/logs/test7/v202608161007/test7\_report.tex}
\par\noindent{\scriptsize TeX 正文；规模：50 行；字节：1\,736；SHA-256：\texttt{a463c7e2b3f105f8}\allowbreak\texttt{944a6459aeb88bc5}\allowbreak\texttt{4013ee281973f12b}\allowbreak\texttt{61e091880475fda0}。}
\item[\texttt{032}] \sourcepath{docs/pyqcd\_reference/from-logs/logs/test8/v202608161117/test8\_report.pdf}
\par\noindent{\scriptsize PDF 见证；规模：5 页；字节：439\,765；SHA-256：\texttt{b7aeac1d21c5d047}\allowbreak\texttt{67fe1580d9473996}\allowbreak\texttt{3add55bd1b0b8304}\allowbreak\texttt{fbce8e8679d330b9}。}
\item[\texttt{033}] \sourcepath{docs/pyqcd\_reference/from-logs/logs/test8/v202608161117/test8\_report.tex}
\par\noindent{\scriptsize TeX 正文；规模：71 行；字节：2\,724；SHA-256：\texttt{ea0f78463994d769}\allowbreak\texttt{d151528811d833fb}\allowbreak\texttt{acd922e249054fcc}\allowbreak\texttt{94b68c5dbba1f115}。}
\item[\texttt{034}] \sourcepath{docs/pyqcd\_reference/from-logs/logs/test9/test9\_analysis.pdf}
\par\noindent{\scriptsize PDF 见证；规模：8 页；字节：692\,938；SHA-256：\texttt{58fba31e5ef000eb}\allowbreak\texttt{a8eebac8e71e7bc9}\allowbreak\texttt{bdd35a4cac8a9f7f}\allowbreak\texttt{e325ca16acc4959a}。}
\item[\texttt{035}] \sourcepath{docs/pyqcd\_reference/from-logs/logs/test9/test9\_analysis.tex}
\par\noindent{\scriptsize TeX 正文；规模：233 行；字节：11\,279；SHA-256：\texttt{74cdf8723a9d7f39}\allowbreak\texttt{45253c0c439b0b96}\allowbreak\texttt{7c56021ad65b7abd}\allowbreak\texttt{909da7a1b439b710}。}
\item[\texttt{036}] \sourcepath{docs/pyqcd\_reference/gevp\_convergence.pdf}
\par\noindent{\scriptsize 独立 PDF；规模：1 页；字节：17\,355；SHA-256：\texttt{32940b99cbd21df5}\allowbreak\texttt{8aa0d34a788accca}\allowbreak\texttt{4c44e8518f1ea83b}\allowbreak\texttt{04dcab5f3e5729c2}。}
\item[\texttt{037}] \sourcepath{docs/pyqcd\_reference/gevp\_t0\_suppression.pdf}
\par\noindent{\scriptsize 独立 PDF；规模：1 页；字节：19\,451；SHA-256：\texttt{3405d3166c006178}\allowbreak\texttt{561532d2dc3c4f5d}\allowbreak\texttt{5b686a140618402f}\allowbreak\texttt{190614409a7ee609}。}
\item[\texttt{038}] \sourcepath{docs/pyqcd\_reference/gradient\_flow.md}
\par\noindent{\scriptsize Markdown；规模：895 行；字节：32\,498；SHA-256：\texttt{b7d8a3e7dc42cd30}\allowbreak\texttt{e89d07c590a19ab0}\allowbreak\texttt{c6063cfb26addb61}\allowbreak\texttt{0dc853f1c8f8a506}。}
\item[\texttt{039}] \sourcepath{docs/pyqcd\_reference/pure\_pyqcd\_20260818.pdf}
\par\noindent{\scriptsize PDF 见证；规模：10 页；字节：340\,582；SHA-256：\texttt{2c0ec28e4a5ec5bb}\allowbreak\texttt{4522897d4d875dc4}\allowbreak\texttt{a0b53bf4698e29d9}\allowbreak\texttt{1055c2c54e5e9d08}。}
\item[\texttt{040}] \sourcepath{docs/pyqcd\_reference/pure\_pyqcd\_20260818.tex}
\par\noindent{\scriptsize TeX 正文；规模：304 行；字节：17\,527；SHA-256：\texttt{cd538e7a7dedd8bb}\allowbreak\texttt{e7a79e57f0ef6908}\allowbreak\texttt{9b21e91f2a1552d5}\allowbreak\texttt{256a53cec084dce9}。}
\item[\texttt{041}] \sourcepath{docs/pyqcd\_reference/report\_cmp1\_4150\_20260828.pdf}
\par\noindent{\scriptsize PDF 见证；规模：16 页；字节：330\,758；SHA-256：\texttt{e03b57a097cca751}\allowbreak\texttt{bb56b50f5e52f687}\allowbreak\texttt{0517a9e4eea3dee3}\allowbreak\texttt{32f3209eeda46fff}。}
\item[\texttt{042}] \sourcepath{docs/pyqcd\_reference/report\_cmp1\_4150\_20260828.tex}
\par\noindent{\scriptsize TeX 正文；规模：507 行；字节：23\,921；SHA-256：\texttt{1cbc3d1c5af60faa}\allowbreak\texttt{03f58139bdb4669e}\allowbreak\texttt{3c73ab1ff165ec50}\allowbreak\texttt{d70922f01fd4af4f}。}
\item[\texttt{043}] \sourcepath{docs/pyqcd\_reference/report\_donghx\_4150\_reproduction\_20260829.pdf}
\par\noindent{\scriptsize PDF 见证；规模：23 页；字节：362\,452；SHA-256：\texttt{ef8d98a11e3e7c98}\allowbreak\texttt{ec0b39f4918bad78}\allowbreak\texttt{dfe3fa4fae7827cf}\allowbreak\texttt{26b026b57bee8ce1}。}
\item[\texttt{044}] \sourcepath{docs/pyqcd\_reference/report\_donghx\_4150\_reproduction\_20260829.tex}
\par\noindent{\scriptsize TeX 正文；规模：617 行；字节：32\,972；SHA-256：\texttt{e235ca69d223bd7f}\allowbreak\texttt{d67e27081e2ef59d}\allowbreak\texttt{d33225a8ac16e4b9}\allowbreak\texttt{ec7ef24509e0dec3}。}
\item[\texttt{045}] \sourcepath{docs/pyqcd\_reference/report\_donghx\_4150\_reproduction\_20260830.pdf}
\par\noindent{\scriptsize PDF 见证；规模：27 页；字节：391\,299；SHA-256：\texttt{f2a0cd8c2f708f97}\allowbreak\texttt{e8986cd07c975968}\allowbreak\texttt{3632f8a0573924e6}\allowbreak\texttt{b23cbc659d3ce4e4}。}
\item[\texttt{046}] \sourcepath{docs/pyqcd\_reference/report\_donghx\_4150\_reproduction\_20260830.tex}
\par\noindent{\scriptsize TeX 正文；规模：732 行；字节：39\,736；SHA-256：\texttt{39117f534b551b71}\allowbreak\texttt{1442cd45070fa354}\allowbreak\texttt{7b68be308c68f0e2}\allowbreak\texttt{d69ffba29611c9a7}。}
\item[\texttt{047}] \sourcepath{docs/pyqcd\_reference/report\_donghx\_pyqcd\_compare\_20260830.pdf}
\par\noindent{\scriptsize PDF 见证；规模：12 页；字节：256\,862；SHA-256：\texttt{8a93ababfe9de0de}\allowbreak\texttt{c1aa16d63c3a6543}\allowbreak\texttt{4045368880996148}\allowbreak\texttt{9414f3300f96ea7d}。}
\item[\texttt{048}] \sourcepath{docs/pyqcd\_reference/report\_donghx\_pyqcd\_compare\_20260830.tex}
\par\noindent{\scriptsize TeX 正文；规模：361 行；字节：19\,419；SHA-256：\texttt{33d203ebe5314e32}\allowbreak\texttt{7975903226a0c62b}\allowbreak\texttt{f44b9af6c0962815}\allowbreak\texttt{03677f15215e2c48}。}
\item[\texttt{049}] \sourcepath{docs/pyqcd\_reference/report\_donghx\_pyqcd\_compare\_20260830\_2.pdf}
\par\noindent{\scriptsize PDF 见证；规模：10 页；字节：151\,746；SHA-256：\texttt{ad1b560a8be2bf97}\allowbreak\texttt{7468a42b6cd9baa2}\allowbreak\texttt{39064d22c9fcc9a2}\allowbreak\texttt{eb8921b646fde207}。}
\item[\texttt{050}] \sourcepath{docs/pyqcd\_reference/report\_donghx\_pyqcd\_compare\_20260830\_2.tex}
\par\noindent{\scriptsize TeX 正文；规模：351 行；字节：20\,056；SHA-256：\texttt{789ceb3160b1a881}\allowbreak\texttt{d6a4ff8313cc50c7}\allowbreak\texttt{f037385ded8ca1ce}\allowbreak\texttt{423714b9671b873f}。}
\item[\texttt{051}] \sourcepath{docs/pyqcd\_reference/report\_gluon\_tmd\_gradient\_flow\_20260828.pdf}
\par\noindent{\scriptsize PDF 见证；规模：22 页；字节：657\,830；SHA-256：\texttt{47038e186fce679f}\allowbreak\texttt{f7b2a098f24ffffc}\allowbreak\texttt{3b4fa34ce303afab}\allowbreak\texttt{d8d4a5672c525c87}。}
\item[\texttt{052}] \sourcepath{docs/pyqcd\_reference/report\_gluon\_tmd\_gradient\_flow\_20260828.tex}
\par\noindent{\scriptsize TeX 正文；规模：789 行；字节：40\,334；SHA-256：\texttt{c8201a1fff3d9bb4}\allowbreak\texttt{54af0c3776b2ad6b}\allowbreak\texttt{e53b8f9f0d5dd312}\allowbreak\texttt{c54209d56fac0713}。}
\item[\texttt{053}] \sourcepath{docs/pyqcd\_reference/report\_gradient\_flow\_gluonpdf\_pyqcd\_20260909.pdf}
\par\noindent{\scriptsize PDF 见证；规模：17 页；字节：450\,305；SHA-256：\texttt{82507699c084eed4}\allowbreak\texttt{588834550cbd2335}\allowbreak\texttt{e586fbfc5c4ab97f}\allowbreak\texttt{114ab490f155dc13}。}
\item[\texttt{054}] \sourcepath{docs/pyqcd\_reference/report\_gradient\_flow\_gluonpdf\_pyqcd\_20260909.tex}
\par\noindent{\scriptsize TeX 正文；规模：1029 行；字节：37\,255；SHA-256：\texttt{a05bc2ba629d23da}\allowbreak\texttt{750264fc5d970aa2}\allowbreak\texttt{9e8a99332d8b1302}\allowbreak\texttt{2a51e69ed6f33ff8}。}
\item[\texttt{055}] \sourcepath{docs/pyqcd\_reference/基于Gradient\_Flow的质子胶子PDF格点计算整体流程.md}
\par\noindent{\scriptsize Markdown；规模：661 行；字节：26\,650；SHA-256：\texttt{a5c1ac4cbe14a543}\allowbreak\texttt{78c2803e9340d4d0}\allowbreak\texttt{704d6ee05f2a4196}\allowbreak\texttt{f7ce644ed6470924}。}
\item[\texttt{056}] \sourcepath{docs/pyqcd\_reference/基于格点QCD对部分子分布函数的研究.pdf}
\par\noindent{\scriptsize 独立 PDF；规模：56 页；字节：2\,466\,379；SHA-256：\texttt{cac44eac6456c007}\allowbreak\texttt{77cc5c4ccaee70a8}\allowbreak\texttt{aa63cbe467fec51e}\allowbreak\texttt{bebb5653fbbd7941}。}
\item[\texttt{057}] \sourcepath{docs/pyqcd\_reference/基于格点QCD的横向动量依赖的部分子分布函数的研究.pdf}
\par\noindent{\scriptsize 独立 PDF；规模：82 页；字节：9\,783\,095；SHA-256：\texttt{dac0a88df96c0d09}\allowbreak\texttt{03e903b319032ca7}\allowbreak\texttt{a270277fb44ef6ac}\allowbreak\texttt{3373813a1b28d2ea}。}
\item[\texttt{058}] \sourcepath{docs/pyqcd\_reference/构造胶子准算符.pdf}
\par\noindent{\scriptsize PDF 见证；规模：14 页；字节：391\,990；SHA-256：\texttt{fb02814754486d5e}\allowbreak\texttt{ea54f70bc88e5d27}\allowbreak\texttt{64563f12f0990427}\allowbreak\texttt{16c6341f00108f30}。}
\item[\texttt{059}] \sourcepath{docs/pyqcd\_reference/构造胶子准算符.tex}
\par\noindent{\scriptsize TeX 正文；规模：562 行；字节：31\,624；SHA-256：\texttt{a160f6f87cba567f}\allowbreak\texttt{f988b7f795477d64}\allowbreak\texttt{68d1bd4934d9bcf7}\allowbreak\texttt{a3b89d37eac51de6}。}
\item[\texttt{060}] \sourcepath{docs/pyqcd\_reference/格点QCD中的3pt构造.pdf}
\par\noindent{\scriptsize PDF 见证；规模：20 页；字节：619\,134；SHA-256：\texttt{d80a8423eb44d646}\allowbreak\texttt{b9f7a066ba77770d}\allowbreak\texttt{d7adf48cd983fdfb}\allowbreak\texttt{c9d2949d5a730997}。}
\item[\texttt{061}] \sourcepath{docs/pyqcd\_reference/格点QCD中的3pt构造.tex}
\par\noindent{\scriptsize TeX 正文；规模：783 行；字节：50\,957；SHA-256：\texttt{573884a330024d4e}\allowbreak\texttt{cafb8d8011a2f2c2}\allowbreak\texttt{3b1f4baef6b616f8}\allowbreak\texttt{e001a5d37306331a}。}
\item[\texttt{062}] \sourcepath{docs/pyqcd\_reference/格点QCD中的DGLAP演化方程.pdf}
\par\noindent{\scriptsize PDF 见证；规模：18 页；字节：528\,085；SHA-256：\texttt{e2c0bdcebf5eaffe}\allowbreak\texttt{b05064734235516e}\allowbreak\texttt{42ea3388847d7337}\allowbreak\texttt{eebf5247cfc2198a}。}
\item[\texttt{063}] \sourcepath{docs/pyqcd\_reference/格点QCD中的DGLAP演化方程.tex}
\par\noindent{\scriptsize TeX 正文；规模：688 行；字节：44\,598；SHA-256：\texttt{a43ca9182e9820bd}\allowbreak\texttt{b8a30558661daa9c}\allowbreak\texttt{26903bd516267f5d}\allowbreak\texttt{b442c554ada8d05e}。}
\item[\texttt{064}] \sourcepath{docs/pyqcd\_reference/格点QCD中的GEVP.pdf}
\par\noindent{\scriptsize PDF 见证；规模：24 页；字节：700\,272；SHA-256：\texttt{ab8efd3a047856fe}\allowbreak\texttt{6a04899417b33739}\allowbreak\texttt{cfdddd57c67a958e}\allowbreak\texttt{a7e4afbac6054c45}。}
\item[\texttt{065}] \sourcepath{docs/pyqcd\_reference/格点QCD中的GEVP.tex}
\par\noindent{\scriptsize TeX 正文；规模：924 行；字节：63\,489；SHA-256：\texttt{6de83102f1088059}\allowbreak\texttt{2f3bc75266dcc2d4}\allowbreak\texttt{8dae6c8ba1e8e4aa}\allowbreak\texttt{8cb8d53047b54d35}。}
\item[\texttt{066}] \sourcepath{docs/pyqcd\_reference/格点QCD中的OPE算符.pdf}
\par\noindent{\scriptsize PDF 见证；规模：14 页；字节：439\,579；SHA-256：\texttt{5c2f643305d95005}\allowbreak\texttt{41489e1e1353a1ac}\allowbreak\texttt{80ea9a48fd2b9018}\allowbreak\texttt{c28ee3c700579271}。}
\item[\texttt{067}] \sourcepath{docs/pyqcd\_reference/格点QCD中的OPE算符.tex}
\par\noindent{\scriptsize TeX 正文；规模：526 行；字节：32\,875；SHA-256：\texttt{3f80d0732fe8e181}\allowbreak\texttt{44b1f3a497f43d75}\allowbreak\texttt{dec7891cc6ac2092}\allowbreak\texttt{3a2f1eb7d094d56b}。}
\item[\texttt{068}] \sourcepath{docs/pyqcd\_reference/格点QCD中的Symanzik有效理论.pdf}
\par\noindent{\scriptsize PDF 见证；规模：13 页；字节：415\,796；SHA-256：\texttt{732f087cd6f3ff57}\allowbreak\texttt{e9a726c0a3b6ba6a}\allowbreak\texttt{12c87905d5e67fa3}\allowbreak\texttt{6b8ba79883892d5f}。}
\item[\texttt{069}] \sourcepath{docs/pyqcd\_reference/格点QCD中的Symanzik有效理论.tex}
\par\noindent{\scriptsize TeX 正文；规模：449 行；字节：29\,755；SHA-256：\texttt{f8abdc1c4aa3f7da}\allowbreak\texttt{3960ea55b3af5a27}\allowbreak\texttt{49f5080f45eff859}\allowbreak\texttt{d6573c06c5fd43e5}。}
\item[\texttt{070}] \sourcepath{docs/pyqcd\_reference/格点QCD中的TMD\_PDF.pdf}
\par\noindent{\scriptsize PDF 见证；规模：30 页；字节：645\,661；SHA-256：\texttt{832ab29c59f2fec2}\allowbreak\texttt{e28f57954cb73ae0}\allowbreak\texttt{c80bd6016865b544}\allowbreak\texttt{2b364d3624af4f3d}。}
\item[\texttt{071}] \sourcepath{docs/pyqcd\_reference/格点QCD中的TMD\_PDF.tex}
\par\noindent{\scriptsize TeX 正文；规模：937 行；字节：72\,101；SHA-256：\texttt{1a785c56d9dd6bad}\allowbreak\texttt{e13b6e5b390d0c08}\allowbreak\texttt{ae7fc5c891973cee}\allowbreak\texttt{a1b89a27982f02d0}。}
\item[\texttt{072}] \sourcepath{docs/pyqcd\_reference/格点QCD中的UV涨落.pdf}
\par\noindent{\scriptsize PDF 见证；规模：13 页；字节：438\,110；SHA-256：\texttt{03fdf94337ff02d1}\allowbreak\texttt{8835fdc5d8e62b51}\allowbreak\texttt{3f3e40449adef1f9}\allowbreak\texttt{8e50375250702c5c}。}
\item[\texttt{073}] \sourcepath{docs/pyqcd\_reference/格点QCD中的UV涨落.tex}
\par\noindent{\scriptsize TeX 正文；规模：472 行；字节：30\,337；SHA-256：\texttt{00a7feb7a3ef32a1}\allowbreak\texttt{b5d7f7b3c893df78}\allowbreak\texttt{3dd640adf28ed587}\allowbreak\texttt{6cb5cfa00e0cb7fe}。}
\item[\texttt{074}] \sourcepath{docs/pyqcd\_reference/格点QCD中的Wilson线.pdf}
\par\noindent{\scriptsize PDF 见证；规模：21 页；字节：559\,057；SHA-256：\texttt{ed84ee969e9f1697}\allowbreak\texttt{8f258cd5e3046cb2}\allowbreak\texttt{ba512f2c6b3eb164}\allowbreak\texttt{a3755a88109a83de}。}
\item[\texttt{075}] \sourcepath{docs/pyqcd\_reference/格点QCD中的Wilson线.tex}
\par\noindent{\scriptsize TeX 正文；规模：712 行；字节：51\,580；SHA-256：\texttt{6ebeae97d4d34fe0}\allowbreak\texttt{8b83f694deb96cab}\allowbreak\texttt{240069bf05d0cadc}\allowbreak\texttt{e27e3d9f4c98831f}。}
\item[\texttt{076}] \sourcepath{docs/pyqcd\_reference/格点QCD中的luscher参数化.pdf}
\par\noindent{\scriptsize PDF 见证；规模：17 页；字节：564\,536；SHA-256：\texttt{f409cc585cef82b8}\allowbreak\texttt{a7008c981c3d5f85}\allowbreak\texttt{fcf9734774e1c9c9}\allowbreak\texttt{122f3f0738259912}。}
\item[\texttt{077}] \sourcepath{docs/pyqcd\_reference/格点QCD中的luscher参数化.tex}
\par\noindent{\scriptsize TeX 正文；规模：603 行；字节：44\,514；SHA-256：\texttt{cea5b0e909fc4a62}\allowbreak\texttt{3bf298b8114e2aad}\allowbreak\texttt{71abc7cfe2141fbe}\allowbreak\texttt{5fdde5ed2ab69357}。}
\item[\texttt{078}] \sourcepath{docs/pyqcd\_reference/格点QCD中的smear算法.pdf}
\par\noindent{\scriptsize PDF 见证；规模：16 页；字节：455\,906；SHA-256：\texttt{2cbd9b2016fd246d}\allowbreak\texttt{11b3a93ac7420bac}\allowbreak\texttt{84bcc21230f1264b}\allowbreak\texttt{896ec4496c2354fe}。}
\item[\texttt{079}] \sourcepath{docs/pyqcd\_reference/格点QCD中的smear算法.tex}
\par\noindent{\scriptsize TeX 正文；规模：630 行；字节：36\,878；SHA-256：\texttt{30685c00813a42b7}\allowbreak\texttt{d4546db135448a0d}\allowbreak\texttt{326033eb68ed4b8f}\allowbreak\texttt{0b08673d31ae8e15}。}
\item[\texttt{080}] \sourcepath{docs/pyqcd\_reference/格点QCD中的传播子求解.pdf}
\par\noindent{\scriptsize PDF 见证；规模：10 页；字节：433\,173；SHA-256：\texttt{7b9d3cb555be3b6d}\allowbreak\texttt{0e584b39931cd4ff}\allowbreak\texttt{7683a4a136d9ab40}\allowbreak\texttt{b473c69a592e762d}。}
\item[\texttt{081}] \sourcepath{docs/pyqcd\_reference/格点QCD中的传播子求解.tex}
\par\noindent{\scriptsize TeX 正文；规模：269 行；字节：22\,693；SHA-256：\texttt{a4c82b384c16bdd9}\allowbreak\texttt{64dcbe7cb9583157}\allowbreak\texttt{6333aa2c8c3bff8a}\allowbreak\texttt{39020b752125af31}。}
\item[\texttt{082}] \sourcepath{docs/pyqcd\_reference/格点QCD中的光锥PDF、准PDF、赝PDF.pdf}
\par\noindent{\scriptsize PDF 见证；规模：20 页；字节：552\,380；SHA-256：\texttt{fb38862b59eb2d42}\allowbreak\texttt{b3aee9aae7de8530}\allowbreak\texttt{74fbbb30a2c66654}\allowbreak\texttt{9c4067688c4da387}。}
\item[\texttt{083}] \sourcepath{docs/pyqcd\_reference/格点QCD中的光锥PDF、准PDF、赝PDF.tex}
\par\noindent{\scriptsize TeX 正文；规模：705 行；字节：48\,477；SHA-256：\texttt{e45e4fb129745406}\allowbreak\texttt{5c04fd8d71096636}\allowbreak\texttt{c5aa7bb52319042f}\allowbreak\texttt{1cc4d5b3f8cd0d23}。}
\item[\texttt{084}] \sourcepath{docs/pyqcd\_reference/格点QCD中的光锥PDF与quasi\_PDF.pdf}
\par\noindent{\scriptsize PDF 见证；规模：12 页；字节：367\,168；SHA-256：\texttt{5d39e651bb6a7400}\allowbreak\texttt{ca2282032e1aafd3}\allowbreak\texttt{690a61fbd861dc0d}\allowbreak\texttt{0f88668f9f5bae11}。}
\item[\texttt{085}] \sourcepath{docs/pyqcd\_reference/格点QCD中的光锥PDF与quasi\_PDF.tex}
\par\noindent{\scriptsize TeX 正文；规模：413 行；字节：25\,659；SHA-256：\texttt{c6cd58cbda714e4c}\allowbreak\texttt{e0221b12caf92958}\allowbreak\texttt{ad66bb21a07de9f5}\allowbreak\texttt{0d06a64b13a8dd8a}。}
\item[\texttt{086}] \sourcepath{docs/pyqcd\_reference/格点QCD中的光锥与光前.pdf}
\par\noindent{\scriptsize PDF 见证；规模：18 页；字节：426\,135；SHA-256：\texttt{cd59c642241b34d0}\allowbreak\texttt{969540dcf636a10d}\allowbreak\texttt{bc1f473704ab9d31}\allowbreak\texttt{38bc125806ef1303}。}
\item[\texttt{087}] \sourcepath{docs/pyqcd\_reference/格点QCD中的光锥与光前.tex}
\par\noindent{\scriptsize TeX 正文；规模：783 行；字节：40\,923；SHA-256：\texttt{d6a20d7aadeabadf}\allowbreak\texttt{bdb5fb1edbd872bf}\allowbreak\texttt{f78871770d44c889}\allowbreak\texttt{cb808f500954f2e5}。}
\item[\texttt{088}] \sourcepath{docs/pyqcd\_reference/格点QCD中的关联函数.pdf}
\par\noindent{\scriptsize PDF 见证；规模：15 页；字节：449\,097；SHA-256：\texttt{5fe68e8c29cd0e61}\allowbreak\texttt{43b31f5046158f3f}\allowbreak\texttt{f35cd40058e2a7e3}\allowbreak\texttt{84e127f07d361231}。}
\item[\texttt{089}] \sourcepath{docs/pyqcd\_reference/格点QCD中的关联函数.tex}
\par\noindent{\scriptsize TeX 正文；规模：604 行；字节：34\,896；SHA-256：\texttt{b761f4a05ff3a6c2}\allowbreak\texttt{0284f0841ad63849}\allowbreak\texttt{8701319ecc52bd36}\allowbreak\texttt{d210e31f30df0ac9}。}
\item[\texttt{090}] \sourcepath{docs/pyqcd\_reference/格点QCD中的场强张量.pdf}
\par\noindent{\scriptsize PDF 见证；规模：14 页；字节：382\,178；SHA-256：\texttt{34ea3d20c45be2e1}\allowbreak\texttt{5717bcb6d4b09d48}\allowbreak\texttt{5eb07dc21d32deeb}\allowbreak\texttt{2e17acffb3d225a7}。}
\item[\texttt{091}] \sourcepath{docs/pyqcd\_reference/格点QCD中的场强张量.tex}
\par\noindent{\scriptsize TeX 正文；规模：576 行；字节：31\,156；SHA-256：\texttt{51c5846c5c8345bf}\allowbreak\texttt{d62d528596b856c8}\allowbreak\texttt{a45c16d0a973efb8}\allowbreak\texttt{8f2e79c18a444fc9}。}
\item[\texttt{092}] \sourcepath{docs/pyqcd\_reference/格点QCD中的外推.pdf}
\par\noindent{\scriptsize PDF 见证；规模：16 页；字节：491\,115；SHA-256：\texttt{c7968256a79eb20e}\allowbreak\texttt{d5b45dcc10c043b5}\allowbreak\texttt{0b54675a2780bb4f}\allowbreak\texttt{e369324bbfb2ccee}。}
\item[\texttt{093}] \sourcepath{docs/pyqcd\_reference/格点QCD中的外推.tex}
\par\noindent{\scriptsize TeX 正文；规模：567 行；字节：40\,966；SHA-256：\texttt{7e0d0e54965cf069}\allowbreak\texttt{c6c51346d8a617d3}\allowbreak\texttt{e1a91af1ef1e5817}\allowbreak\texttt{d4a8ca59ead7b443}。}
\item[\texttt{094}] \sourcepath{docs/pyqcd\_reference/格点QCD中的大动量有效理论.pdf}
\par\noindent{\scriptsize PDF 见证；规模：13 页；字节：415\,295；SHA-256：\texttt{a07c6a93b3c2c71a}\allowbreak\texttt{a1e312be0d5baeb0}\allowbreak\texttt{56ec08c00e421291}\allowbreak\texttt{9d078d389b241226}。}
\item[\texttt{095}] \sourcepath{docs/pyqcd\_reference/格点QCD中的大动量有效理论.tex}
\par\noindent{\scriptsize TeX 正文；规模：498 行；字节：29\,104；SHA-256：\texttt{ece25f53cba99342}\allowbreak\texttt{9fe85c19f418a750}\allowbreak\texttt{e66e0b982a800069}\allowbreak\texttt{309e9bb1ff4983aa}。}
\item[\texttt{096}] \sourcepath{docs/pyqcd\_reference/格点QCD中的夸克算符.pdf}
\par\noindent{\scriptsize PDF 见证；规模：13 页；字节：402\,811；SHA-256：\texttt{b848bf53dbe2df2c}\allowbreak\texttt{22f12ff8d43ed962}\allowbreak\texttt{c3902c83837ae544}\allowbreak\texttt{40e925dec1038e8b}。}
\item[\texttt{097}] \sourcepath{docs/pyqcd\_reference/格点QCD中的夸克算符.tex}
\par\noindent{\scriptsize TeX 正文；规模：457 行；字节：29\,049；SHA-256：\texttt{fc960e0f9fe18a12}\allowbreak\texttt{3e4523b7f1261b59}\allowbreak\texttt{9eb40d870d24420b}\allowbreak\texttt{122fe0dda2311e08}。}
\item[\texttt{098}] \sourcepath{docs/pyqcd\_reference/格点QCD中的常见名词短语与缩写.pdf}
\par\noindent{\scriptsize PDF 见证；规模：26 页；字节：532\,740；SHA-256：\texttt{23e63f872e382306}\allowbreak\texttt{f1d27756d20e953e}\allowbreak\texttt{4352e39f3c4af9da}\allowbreak\texttt{229eecc51ef607d4}。}
\item[\texttt{099}] \sourcepath{docs/pyqcd\_reference/格点QCD中的常见名词短语与缩写.tex}
\par\noindent{\scriptsize TeX 正文；规模：414 行；字节：55\,807；SHA-256：\texttt{5923a71cd886bf19}\allowbreak\texttt{cfb2e64913178985}\allowbreak\texttt{437df285a997a192}\allowbreak\texttt{887eb2e9606ab5ad}。}
\item[\texttt{100}] \sourcepath{docs/pyqcd\_reference/格点QCD中的常见费曼图.pdf}
\par\noindent{\scriptsize PDF 见证；规模：10 页；字节：349\,511；SHA-256：\texttt{37ca0f40e748ac16}\allowbreak\texttt{bd530576ba24d93d}\allowbreak\texttt{f0e6a2674c6cde1f}\allowbreak\texttt{9d543e42e9966347}。}
\item[\texttt{101}] \sourcepath{docs/pyqcd\_reference/格点QCD中的常见费曼图.tex}
\par\noindent{\scriptsize TeX 正文；规模：495 行；字节：26\,316；SHA-256：\texttt{bfce4234f827be27}\allowbreak\texttt{c4aa1ab9c0fba9c3}\allowbreak\texttt{677d4a7acd664bbe}\allowbreak\texttt{f912a57cbd60eda5}。}
\item[\texttt{102}] \sourcepath{docs/pyqcd\_reference/格点QCD中的广义部分子分布.pdf}
\par\noindent{\scriptsize PDF 见证；规模：25 页；字节：545\,329；SHA-256：\texttt{60486936257c8161}\allowbreak\texttt{a209160f5cd473c3}\allowbreak\texttt{070207c60d5cc039}\allowbreak\texttt{b73ab89462e22959}。}
\item[\texttt{103}] \sourcepath{docs/pyqcd\_reference/格点QCD中的广义部分子分布.tex}
\par\noindent{\scriptsize TeX 正文；规模：1354 行；字节：55\,674；SHA-256：\texttt{6436bbca893f8242}\allowbreak\texttt{46b9dd295c237b0d}\allowbreak\texttt{60c6760fcc32dd6f}\allowbreak\texttt{5935821b95c71fd2}。}
\item[\texttt{104}] \sourcepath{docs/pyqcd\_reference/格点QCD中的强子谱学.pdf}
\par\noindent{\scriptsize PDF 见证；规模：11 页；字节：439\,197；SHA-256：\texttt{b8872bb3cf9fabaa}\allowbreak\texttt{9ae5a80e43a9949d}\allowbreak\texttt{646376fb1d6388a9}\allowbreak\texttt{02f41e865146d1dd}。}
\item[\texttt{105}] \sourcepath{docs/pyqcd\_reference/格点QCD中的强子谱学.tex}
\par\noindent{\scriptsize TeX 正文；规模：353 行；字节：25\,165；SHA-256：\texttt{973ad5dcbd50830e}\allowbreak\texttt{d82df520be6dd76b}\allowbreak\texttt{d3ab64e783255ca3}\allowbreak\texttt{76238dc56f5cb0a3}。}
\item[\texttt{106}] \sourcepath{docs/pyqcd\_reference/格点QCD中的形状因子.pdf}
\par\noindent{\scriptsize PDF 见证；规模：14 页；字节：511\,026；SHA-256：\texttt{91da8f2e52c3ca24}\allowbreak\texttt{3d156f81ffd7955c}\allowbreak\texttt{2386418967d49a63}\allowbreak\texttt{ac9245fe7d2d75ac}。}
\item[\texttt{107}] \sourcepath{docs/pyqcd\_reference/格点QCD中的形状因子.tex}
\par\noindent{\scriptsize TeX 正文；规模：499 行；字节：36\,680；SHA-256：\texttt{4abc0fb2b9904a77}\allowbreak\texttt{7524058f9561bc3e}\allowbreak\texttt{89fe809a1a431b29}\allowbreak\texttt{4d1d325231f82e88}。}
\item[\texttt{108}] \sourcepath{docs/pyqcd\_reference/格点QCD中的标度参数.pdf}
\par\noindent{\scriptsize PDF 见证；规模：13 页；字节：415\,188；SHA-256：\texttt{a1f22b621bd77a14}\allowbreak\texttt{bbfa2554bd6a3e7f}\allowbreak\texttt{bc6a37d2e3269d2c}\allowbreak\texttt{73508b5046c5ceba}。}
\item[\texttt{109}] \sourcepath{docs/pyqcd\_reference/格点QCD中的标度参数.tex}
\par\noindent{\scriptsize TeX 正文；规模：456 行；字节：28\,707；SHA-256：\texttt{3523d6f1ff1fa017}\allowbreak\texttt{3079a6970592544d}\allowbreak\texttt{620989424a8e7130}\allowbreak\texttt{c87aabee537bb6e4}。}
\item[\texttt{110}] \sourcepath{docs/pyqcd\_reference/格点QCD中的梯度流重整化.pdf}
\par\noindent{\scriptsize PDF 见证；规模：15 页；字节：429\,035；SHA-256：\texttt{586a27b15c8e3bd1}\allowbreak\texttt{4330da6dbefd2c55}\allowbreak\texttt{90d49f6d1c1dd0b2}\allowbreak\texttt{11c518e447ad8803}。}
\item[\texttt{111}] \sourcepath{docs/pyqcd\_reference/格点QCD中的梯度流重整化.tex}
\par\noindent{\scriptsize TeX 正文；规模：626 行；字节：37\,208；SHA-256：\texttt{8717b2a859f7deec}\allowbreak\texttt{8714828e4325a730}\allowbreak\texttt{48ed9f71c0d5df5f}\allowbreak\texttt{b66404db56db44a8}。}
\item[\texttt{112}] \sourcepath{docs/pyqcd\_reference/格点QCD中的维克收缩与傅里叶变换.pdf}
\par\noindent{\scriptsize PDF 见证；规模：13 页；字节：425\,972；SHA-256：\texttt{046e5cc683b39afd}\allowbreak\texttt{08523201a018b738}\allowbreak\texttt{e7ecad5247f2c8a5}\allowbreak\texttt{c354a51a31dd4005}。}
\item[\texttt{113}] \sourcepath{docs/pyqcd\_reference/格点QCD中的维克收缩与傅里叶变换.tex}
\par\noindent{\scriptsize TeX 正文；规模：469 行；字节：29\,279；SHA-256：\texttt{28a3d50b1f48d74f}\allowbreak\texttt{b51f23b32c919f6b}\allowbreak\texttt{d8d1af3f7e857a1f}\allowbreak\texttt{621e388f9425709b}。}
\item[\texttt{114}] \sourcepath{docs/pyqcd\_reference/格点QCD中的胶子极化.pdf}
\par\noindent{\scriptsize PDF 见证；规模：16 页；字节：443\,418；SHA-256：\texttt{db5bd54dd07d031d}\allowbreak\texttt{f400d2f2438e88a9}\allowbreak\texttt{ee6e0cc41ef67f2f}\allowbreak\texttt{51d7a22e3792b67b}。}
\item[\texttt{115}] \sourcepath{docs/pyqcd\_reference/格点QCD中的胶子极化.tex}
\par\noindent{\scriptsize TeX 正文；规模：694 行；字节：38\,655；SHA-256：\texttt{1cef9c8e9fae72a1}\allowbreak\texttt{a43620e4e54affe8}\allowbreak\texttt{bc4d51348ef7936b}\allowbreak\texttt{9398d272636752f7}。}
\item[\texttt{116}] \sourcepath{docs/pyqcd\_reference/格点QCD中的胶子算符.pdf}
\par\noindent{\scriptsize PDF 见证；规模：15 页；字节：396\,500；SHA-256：\texttt{970ef249969bd25c}\allowbreak\texttt{3787a03d86105a2e}\allowbreak\texttt{eef736e6e08dcec6}\allowbreak\texttt{db7a2e6a6b183743}。}
\item[\texttt{117}] \sourcepath{docs/pyqcd\_reference/格点QCD中的胶子算符.tex}
\par\noindent{\scriptsize TeX 正文；规模：618 行；字节：34\,339；SHA-256：\texttt{644f1066ef957e02}\allowbreak\texttt{7216845aaa775430}\allowbreak\texttt{9e5b804ba32ff06f}\allowbreak\texttt{82b730674bcaf173}。}
\item[\texttt{118}] \sourcepath{docs/pyqcd\_reference/格点QCD中的色散关系.pdf}
\par\noindent{\scriptsize PDF 见证；规模：15 页；字节：504\,778；SHA-256：\texttt{d03a804df50ea2f2}\allowbreak\texttt{81cc8da217d49511}\allowbreak\texttt{a47a2471f4b60dfc}\allowbreak\texttt{61e99c8a78526456}。}
\item[\texttt{119}] \sourcepath{docs/pyqcd\_reference/格点QCD中的色散关系.tex}
\par\noindent{\scriptsize TeX 正文；规模：496 行；字节：37\,730；SHA-256：\texttt{0737abf85734aa31}\allowbreak\texttt{27bb8b9d4046b0f1}\allowbreak\texttt{5d3f5dd248b6d2f0}\allowbreak\texttt{43a21b7c210a53d6}。}
\item[\texttt{120}] \sourcepath{docs/pyqcd\_reference/格点QCD中的蒙卡方法.pdf}
\par\noindent{\scriptsize PDF 见证；规模：11 页；字节：462\,990；SHA-256：\texttt{444777355609bc8f}\allowbreak\texttt{6fbd17787e08833e}\allowbreak\texttt{0a2052859870f3f0}\allowbreak\texttt{cd3d810aebf0ab65}。}
\item[\texttt{121}] \sourcepath{docs/pyqcd\_reference/格点QCD中的蒙卡方法.tex}
\par\noindent{\scriptsize TeX 正文；规模：347 行；字节：28\,178；SHA-256：\texttt{a33d4f7d9e67cb93}\allowbreak\texttt{a1b2337bf4d1fc31}\allowbreak\texttt{e5c736a582ae6d96}\allowbreak\texttt{fe7a26d418dfc1ee}。}
\item[\texttt{122}] \sourcepath{docs/pyqcd\_reference/格点QCD中的衰变道.pdf}
\par\noindent{\scriptsize PDF 见证；规模：15 页；字节：512\,657；SHA-256：\texttt{373c1dc85d9ae551}\allowbreak\texttt{8be0df30128d0832}\allowbreak\texttt{d48545fa02b33b5a}\allowbreak\texttt{ddd5030f424582b6}。}
\item[\texttt{123}] \sourcepath{docs/pyqcd\_reference/格点QCD中的衰变道.tex}
\par\noindent{\scriptsize TeX 正文；规模：474 行；字节：37\,402；SHA-256：\texttt{a72d61b1ef1c9217}\allowbreak\texttt{bfd58efad717c822}\allowbreak\texttt{10ee05fd19b731e2}\allowbreak\texttt{e6d7ada7634a0aed}。}
\item[\texttt{124}] \sourcepath{docs/pyqcd\_reference/格点QCD中的误差统计.pdf}
\par\noindent{\scriptsize PDF 见证；规模：15 页；字节：490\,086；SHA-256：\texttt{e7d90bb41b7d7145}\allowbreak\texttt{304de93244ca5570}\allowbreak\texttt{813d3885aa73f597}\allowbreak\texttt{66d28f4e0edb8f2a}。}
\item[\texttt{125}] \sourcepath{docs/pyqcd\_reference/格点QCD中的误差统计.tex}
\par\noindent{\scriptsize TeX 正文；规模：472 行；字节：36\,291；SHA-256：\texttt{017f3acd57cdc449}\allowbreak\texttt{6ea24c08da1bf620}\allowbreak\texttt{028aae08fb2a3a42}\allowbreak\texttt{4f2d9423dd14422f}。}
\item[\texttt{126}] \sourcepath{docs/pyqcd\_reference/格点QCD中的费米子方案.pdf}
\par\noindent{\scriptsize PDF 见证；规模：24 页；字节：576\,602；SHA-256：\texttt{832e3f7f46ea0f83}\allowbreak\texttt{33bb4211d2a60053}\allowbreak\texttt{ebe537ea0a8c836e}\allowbreak\texttt{a84b731e9736e376}。}
\item[\texttt{127}] \sourcepath{docs/pyqcd\_reference/格点QCD中的费米子方案.tex}
\par\noindent{\scriptsize TeX 正文；规模：999 行；字节：61\,284；SHA-256：\texttt{a675522d85691d24}\allowbreak\texttt{5694fcfe6ddf5e82}\allowbreak\texttt{2c57aa52c26eb142}\allowbreak\texttt{3a191c77a27f6018}。}
\item[\texttt{128}] \sourcepath{docs/pyqcd\_reference/格点QCD中的连通图与非连通图.pdf}
\par\noindent{\scriptsize PDF 见证；规模：16 页；字节：537\,089；SHA-256：\texttt{46cb1b5de3a56184}\allowbreak\texttt{e78f635dbba24a30}\allowbreak\texttt{96a5fe1e0b7921f2}\allowbreak\texttt{72bfcb947f1d2158}。}
\item[\texttt{129}] \sourcepath{docs/pyqcd\_reference/格点QCD中的连通图与非连通图.tex}
\par\noindent{\scriptsize TeX 正文；规模：531 行；字节：42\,019；SHA-256：\texttt{4120d271fe16fdb9}\allowbreak\texttt{3d2deb344f93d908}\allowbreak\texttt{724b1070b0edc56c}\allowbreak\texttt{c5322b0585a94fa8}。}
\item[\texttt{130}] \sourcepath{docs/pyqcd\_reference/格点QCD中的部分子分布函数.pdf}
\par\noindent{\scriptsize PDF 见证；规模：16 页；字节：454\,948；SHA-256：\texttt{ec25264bf8ed9910}\allowbreak\texttt{72f1618c6870104f}\allowbreak\texttt{696c0c78df0f5904}\allowbreak\texttt{4a911c7f5a13746f}。}
\item[\texttt{131}] \sourcepath{docs/pyqcd\_reference/格点QCD中的部分子分布函数.tex}
\par\noindent{\scriptsize TeX 正文；规模：576 行；字节：34\,642；SHA-256：\texttt{ab83c2e090ab9639}\allowbreak\texttt{053ed7c0d75156d0}\allowbreak\texttt{3433a7328ca5bc73}\allowbreak\texttt{72e4b35493996091}。}
\item[\texttt{132}] \sourcepath{docs/pyqcd\_reference/格点QCD中的重夸克有效理论.pdf}
\par\noindent{\scriptsize PDF 见证；规模：12 页；字节：420\,832；SHA-256：\texttt{88682031dcc17573}\allowbreak\texttt{6c37133975604316}\allowbreak\texttt{cbfb658710595177}\allowbreak\texttt{5463c65d478515a5}。}
\item[\texttt{133}] \sourcepath{docs/pyqcd\_reference/格点QCD中的重夸克有效理论.tex}
\par\noindent{\scriptsize TeX 正文；规模：426 行；字节：27\,156；SHA-256：\texttt{7a7e801e6a3ae90a}\allowbreak\texttt{b7a1b0a199cdee71}\allowbreak\texttt{d961eac22a1d2581}\allowbreak\texttt{8e86d4aaf4bbff83}。}
\item[\texttt{134}] \sourcepath{docs/pyqcd\_reference/格点QCD中的重整化.pdf}
\par\noindent{\scriptsize PDF 见证；规模：21 页；字节：567\,098；SHA-256：\texttt{2ad98239dc1d640d}\allowbreak\texttt{86de9b1080f929d0}\allowbreak\texttt{53f05191dac715bb}\allowbreak\texttt{c90b09d058d27251}。}
\item[\texttt{135}] \sourcepath{docs/pyqcd\_reference/格点QCD中的重整化.tex}
\par\noindent{\scriptsize TeX 正文；规模：741 行；字节：52\,997；SHA-256：\texttt{59bb49dd32d10e4f}\allowbreak\texttt{9ebb580e7edf8be9}\allowbreak\texttt{ad7248f4379d7263}\allowbreak\texttt{6420d87f0f0b01da}。}
\item[\texttt{136}] \sourcepath{docs/pyqcd\_reference/格点QCD中的重采样方法.pdf}
\par\noindent{\scriptsize PDF 见证；规模：16 页；字节：517\,180；SHA-256：\texttt{aa90d94361af845c}\allowbreak\texttt{a2b0831eaf56e294}\allowbreak\texttt{620b311ace5a332b}\allowbreak\texttt{3a42f843153da3cf}。}
\item[\texttt{137}] \sourcepath{docs/pyqcd\_reference/格点QCD中的重采样方法.tex}
\par\noindent{\scriptsize TeX 正文；规模：507 行；字节：40\,132；SHA-256：\texttt{2dbc83ca816afaf4}\allowbreak\texttt{eea0e062f08bb583}\allowbreak\texttt{357c3b34599a8e64}\allowbreak\texttt{e054f653668be336}。}
\item[\texttt{138}] \sourcepath{docs/pyqcd\_reference/格点QCD蒸馏方法解析.pdf}
\par\noindent{\scriptsize PDF 见证；规模：14 页；字节：411\,181；SHA-256：\texttt{c1d3d4e36e46778b}\allowbreak\texttt{fe41036e73433f7d}\allowbreak\texttt{c448da662aa917a8}\allowbreak\texttt{5f8f52b095e9cdde}。}
\item[\texttt{139}] \sourcepath{docs/pyqcd\_reference/格点QCD蒸馏方法解析.tex}
\par\noindent{\scriptsize TeX 正文；规模：586 行；字节：31\,130；SHA-256：\texttt{ffe1ef95188fa8c6}\allowbreak\texttt{d2576b7a554046e7}\allowbreak\texttt{c14dfb0a455bff0c}\allowbreak\texttt{7faec4207db584b7}。}
\item[\texttt{140}] \sourcepath{docs/pyqcd\_reference/格点上计算胶子准算符.pdf}
\par\noindent{\scriptsize PDF 见证；规模：22 页；字节：555\,215；SHA-256：\texttt{a3bf39abebee5cfc}\allowbreak\texttt{66a5c211aca4f434}\allowbreak\texttt{97155e59ca985c8f}\allowbreak\texttt{4ed0c83c16f6baeb}。}
\item[\texttt{141}] \sourcepath{docs/pyqcd\_reference/格点上计算胶子准算符.tex}
\par\noindent{\scriptsize TeX 正文；规模：947 行；字节：49\,800；SHA-256：\texttt{4a141f1dd1b20aaf}\allowbreak\texttt{eeb73c94dd2467ad}\allowbreak\texttt{25da919f1b8bd900}\allowbreak\texttt{dd07c8785c1bac5f}。}
\item[\texttt{142}] \sourcepath{docs/pyqcd\_reference/理论解析与工作流-v20260911.pdf}
\par\noindent{\scriptsize PDF 见证；规模：100 页；字节：1\,373\,004；SHA-256：\texttt{ee275800365b5193}\allowbreak\texttt{0539d3a742f9312e}\allowbreak\texttt{4220c33841a42831}\allowbreak\texttt{3e0b8b141b16b34a}。}
\item[\texttt{143}] \sourcepath{docs/pyqcd\_reference/理论解析与工作流-v20260911.tex}
\par\noindent{\scriptsize TeX 正文；规模：5957 行；字节：300\,841；SHA-256：\texttt{699e9f3060838952}\allowbreak\texttt{1b0d2f82a8047f2d}\allowbreak\texttt{08ef9ed40bc1a60d}\allowbreak\texttt{9715e44a7796f1d3}。}
\item[\texttt{144}] \sourcepath{docs/pyqcd\_reference/理论解析与工作流.pdf}
\par\noindent{\scriptsize PDF 见证；规模：89 页；字节：1\,142\,493；SHA-256：\texttt{16601e44ccd3e852}\allowbreak\texttt{dc2763e405a8bbae}\allowbreak\texttt{a9e3898583dba10c}\allowbreak\texttt{d5cc7e0a864e8925}。}
\item[\texttt{145}] \sourcepath{docs/pyqcd\_reference/理论解析与工作流.tex}
\par\noindent{\scriptsize TeX 正文；规模：3837 行；字节：220\,708；SHA-256：\texttt{c9e5e0d25c50bee4}\allowbreak\texttt{f4c5220e380e0fe6}\allowbreak\texttt{835b947e771e35ab}\allowbreak\texttt{8d17046245963e0c}。}
\item[\texttt{146}] \sourcepath{docs/pyqcd\_reference/胶子PDF完整推导.pdf}
\par\noindent{\scriptsize PDF 见证；规模：15 页；字节：348\,290；SHA-256：\texttt{9e660d0139702b1c}\allowbreak\texttt{77b0b907987fd926}\allowbreak\texttt{a7f5b2b9cc078218}\allowbreak\texttt{0019d7d9bee22d77}。}
\item[\texttt{147}] \sourcepath{docs/pyqcd\_reference/胶子PDF完整推导.tex}
\par\noindent{\scriptsize TeX 正文；规模：742 行；字节：37\,897；SHA-256：\texttt{087e590d5818e495}\allowbreak\texttt{f0ac08aa5931bbd4}\allowbreak\texttt{e9bf2a5cb0436efa}\allowbreak\texttt{aef608e2094e355e}。}
\item[\texttt{148}] \sourcepath{docs/pyqcd\_reference/胶子PDF幻灯片.pdf}
\par\noindent{\scriptsize PDF 见证；规模：10 页；字节：219\,085；SHA-256：\texttt{1eb005812e1fd9e1}\allowbreak\texttt{233f9cabb0491251}\allowbreak\texttt{682f47575be78c2d}\allowbreak\texttt{830775c89947eed9}。}
\item[\texttt{149}] \sourcepath{docs/pyqcd\_reference/胶子PDF幻灯片.tex}
\par\noindent{\scriptsize TeX 正文；规模：357 行；字节：16\,128；SHA-256：\texttt{d2adf3c955782d99}\allowbreak\texttt{1ca1b142757d4308}\allowbreak\texttt{1d5e9d216cd7c799}\allowbreak\texttt{ac2b2dbb8dea7a6e}。}
\item[\texttt{150}] \sourcepath{docs/pyqcd\_reference/胶子PDF笔记.pdf}
\par\noindent{\scriptsize PDF 见证；规模：9 页；字节：283\,294；SHA-256：\texttt{7504e1c11dfd3412}\allowbreak\texttt{86a1333b572713ba}\allowbreak\texttt{02e03cf58ccac829}\allowbreak\texttt{9a3e24543574ebb8}。}
\item[\texttt{151}] \sourcepath{docs/pyqcd\_reference/胶子PDF笔记.tex}
\par\noindent{\scriptsize TeX 正文；规模：445 行；字节：22\,729；SHA-256：\texttt{f96be544e01665f9}\allowbreak\texttt{90019414e61d713e}\allowbreak\texttt{e587496fb1359a04}\allowbreak\texttt{ec718d7b52c04612}。}
\item[\texttt{152}] \sourcepath{docs/pyqcd\_reference/胶子PDF计算代码解析.pdf}
\par\noindent{\scriptsize PDF 见证；规模：30 页；字节：462\,296；SHA-256：\texttt{2ad33f213b61daae}\allowbreak\texttt{ba5426c428d58c58}\allowbreak\texttt{3a08e54e42559846}\allowbreak\texttt{8336bb945115c479}。}
\item[\texttt{153}] \sourcepath{docs/pyqcd\_reference/胶子PDF计算代码解析.tex}
\par\noindent{\scriptsize TeX 正文；规模：1218 行；字节：51\,688；SHA-256：\texttt{d12460de77c87586}\allowbreak\texttt{7bd2bad98eb1ec18}\allowbreak\texttt{19f4c8633ffc3aaf}\allowbreak\texttt{5c4e22d69a23151d}。}
\item[\texttt{154}] \sourcepath{docs/pyqcd\_reference/胶子PDF连续极限.pdf}
\par\noindent{\scriptsize PDF 见证；规模：16 页；字节：372\,848；SHA-256：\texttt{2061cd46250653c3}\allowbreak\texttt{843c51bb769d817c}\allowbreak\texttt{6401799a19ebe83e}\allowbreak\texttt{14354121402e05a3}。}
\item[\texttt{155}] \sourcepath{docs/pyqcd\_reference/胶子PDF连续极限.tex}
\par\noindent{\scriptsize TeX 正文；规模：578 行；字节：37\,759；SHA-256：\texttt{9533deecbd993ee2}\allowbreak\texttt{fc1ebba4bc43c3b4}\allowbreak\texttt{5f44663e3e4176f5}\allowbreak\texttt{e2bc89c6764e4515}。}
\item[\texttt{156}] \sourcepath{docs/pyqcd\_reference/胶子PDF连续极限幻灯片.pdf}
\par\noindent{\scriptsize PDF 见证；规模：50 页；字节：405\,678；SHA-256：\texttt{4c25a83f40752f64}\allowbreak\texttt{740a747ee56b6c3f}\allowbreak\texttt{cd3737732a36d94c}\allowbreak\texttt{23d8f6c8aac0d483}。}
\item[\texttt{157}] \sourcepath{docs/pyqcd\_reference/胶子PDF连续极限幻灯片.tex}
\par\noindent{\scriptsize TeX 正文；规模：1576 行；字节：64\,393；SHA-256：\texttt{1cbb0909185b1ccb}\allowbreak\texttt{fc5f7e6a653db68c}\allowbreak\texttt{515486bdbf1318b1}\allowbreak\texttt{c125c6930d29752b}。}
\item[\texttt{158}] \sourcepath{docs/pyqcd\_reference/胶子PDF连续极限论文.pdf}
\par\noindent{\scriptsize PDF 见证；规模：16 页；字节：371\,424；SHA-256：\texttt{565122646ad7873b}\allowbreak\texttt{50e2dc333f1e399c}\allowbreak\texttt{10bcfe7319e56e75}\allowbreak\texttt{103ea97790333ac6}。}
\item[\texttt{159}] \sourcepath{docs/pyqcd\_reference/胶子PDF连续极限论文.tex}
\par\noindent{\scriptsize TeX 正文；规模：573 行；字节：37\,479；SHA-256：\texttt{86160991a9ef06e3}\allowbreak\texttt{d699f7341ee9a0ad}\allowbreak\texttt{b2473342de2a5f9d}\allowbreak\texttt{695221577f0ecf88}。}
\item[\texttt{160}] \sourcepath{docs/pyqcd\_reference/胶子TMD梯度流连续极限.pdf}
\par\noindent{\scriptsize PDF 见证；规模：35 页；字节：585\,221；SHA-256：\texttt{fc507ef05992916e}\allowbreak\texttt{db6249977e702e76}\allowbreak\texttt{247119ed5fb2d0c6}\allowbreak\texttt{20c1e6dafcf443ae}。}
\item[\texttt{161}] \sourcepath{docs/pyqcd\_reference/胶子TMD梯度流连续极限.tex}
\par\noindent{\scriptsize TeX 正文；规模：1481 行；字节：77\,181；SHA-256：\texttt{34f56b3b909edd95}\allowbreak\texttt{64e4aa37416f97dd}\allowbreak\texttt{7d40b090c2b398b1}\allowbreak\texttt{e160951d0798f726}。}
\item[\texttt{162}] \sourcepath{docs/pyqcd\_reference/质子自旋危机解析.pdf}
\par\noindent{\scriptsize PDF 见证；规模：14 页；字节：416\,167；SHA-256：\texttt{229f03499407ad3e}\allowbreak\texttt{fb65322fe163cc8e}\allowbreak\texttt{5a651d961481f1b9}\allowbreak\texttt{8f6d64da42a35612}。}
\item[\texttt{163}] \sourcepath{docs/pyqcd\_reference/质子自旋危机解析.tex}
\par\noindent{\scriptsize TeX 正文；规模：564 行；字节：30\,901；SHA-256：\texttt{ba016d3810d81819}\allowbreak\texttt{1a41cfe3f1ea7dc2}\allowbreak\texttt{f8bfe5de09142feb}\allowbreak\texttt{7c8b426168be22b6}。}
\end{description}
\section{控制文件}
以下文件不进入物理内容计数，但决定快照边界和许可证解释。
\begin{description}[leftmargin=4.2em,labelwidth=3.7em,style=nextline,itemsep=0.25em]
\item[\texttt{C01}] \sourcepath{docs/pyqcd\_reference/AGENTS.md}
\par\noindent{\scriptsize 控制文件；规模：14 行；字节：857；SHA-256：\texttt{896c08fb62b3b348}\allowbreak\texttt{432bbc404da427e3}\allowbreak\texttt{50301de9724c2f2f}\allowbreak\texttt{4b5cf1c4682298b1}。}
\item[\texttt{C02}] \sourcepath{docs/pyqcd\_reference/LICENSE.md}
\par\noindent{\scriptsize 控制文件；规模：21 行；字节：1\,071；SHA-256：\texttt{edf5f3be1ced1214}\allowbreak\texttt{81c53c774c321a95}\allowbreak\texttt{5c8c6eca1f188f1e}\allowbreak\texttt{6cdeac0d00216776}。}
\item[\texttt{C03}] \sourcepath{docs/pyqcd\_reference/README.md}
\par\noindent{\scriptsize 控制文件；规模：60 行；字节：3\,243；SHA-256：\texttt{71a3f3e53ed91608}\allowbreak\texttt{de8fa72412cd712f}\allowbreak\texttt{9540b337dd669595}\allowbreak\texttt{7933fa17212a6d7b}。}
\end{description}
\section{覆盖结论与边界}
\begin{checkpointbox}{全量覆盖判据}
\small
\begin{enumerate}
  \item 163 个内容文件均有 SHA-256 和章节归宿；76 对 TeX/PDF 以同一知识单元计，但两个文件均登记。
  \item 历史重复版本不会被当作新物理结论重复堆叠；其差异用于验证演进、失败修复和实现状态变化。
  \item PDF 中的数值结果只按快照语境引用；除非专著另行登记当前源码、运行输入和验证输出，否则不升级为新的物理验收结果。
  \item 固定的上游历史路径、旧接口和集群命令仅属于引用语境，不构成本库当前执行入口。
\end{enumerate}
\end{checkpointbox}
